\documentclass[11pt,a4paper]{article}

\input{glyphtounicode}
\pdfgentounicode=1
\usepackage[utf8]{inputenc}
\usepackage{cmap}
\usepackage[T1]{fontenc}
\usepackage{lmodern}
\usepackage{amsmath,amssymb,amsthm,mathtools}
\usepackage{booktabs}
\usepackage{enumitem}
\usepackage{geometry}
\usepackage{makecell}
\usepackage{xurl}
\usepackage[colorlinks=true,linkcolor=blue,citecolor=blue,urlcolor=blue]{hyperref}
\usepackage[nameinlink,capitalize]{cleveref}

\geometry{margin=1in}
\setlength{\emergencystretch}{2em}

\newtheorem{theorem}{Theorem}[section]
\newtheorem{lemma}[theorem]{Lemma}
\newtheorem{proposition}[theorem]{Proposition}
\newtheorem{corollary}[theorem]{Corollary}
\newtheorem{definition}[theorem]{Definition}
\theoremstyle{remark}
\newtheorem{remark}[theorem]{Remark}
\theoremstyle{plain}
\newtheorem{conjecture}[theorem]{Conjecture}
\crefname{theorem}{Theorem}{Theorems}
\crefname{lemma}{Lemma}{Lemmas}
\crefname{proposition}{Proposition}{Propositions}
\crefname{corollary}{Corollary}{Corollaries}
\crefname{definition}{Definition}{Definitions}
\crefname{remark}{Remark}{Remarks}
\crefname{conjecture}{Conjecture}{Conjectures}

\newcommand{\QW}{\mathrm{QW}}
\newcommand{\RR}{\mathbb{R}}
\newcommand{\PP}{\mathbb{P}}
\newcommand{\norm}[1]{\left\lVert #1 \right\rVert}
\newcommand{\projectrepo}{\href{https://github.com/research-line/rh-even-dominance}{\path{github.com/research-line/rh-even-dominance}}}
\DeclareMathOperator{\Tr}{Tr}

\title{A Conditional Reduction of the Riemann Hypothesis to Even Dominance of the Weil Quadratic Form}
\author{Lukas Geiger\\Independent Researcher, Bernau\\ORCID: 0009-0005-7296-1534}
\date{Draft v1.7 -- 23 May 2026}

\hypersetup{
  unicode=true,
  pdftitle={A Conditional Reduction of the Riemann Hypothesis to Even Dominance of the Weil Quadratic Form},
  pdfauthor={Lukas Geiger},
  pdfsubject={Conditional reduction, diagnostic programme, and numerical evidence toward the Riemann Hypothesis},
  pdfkeywords={Riemann Hypothesis, Weil quadratic form, even dominance, conditional reduction, Connes program}
}

\begin{document}
\maketitle

\begin{abstract}
This paper presents a conditional reduction of the Riemann Hypothesis,
extracted from a three-part Landscape series. The Landscape papers document
the full research process, including failed routes, obstruction theorems,
computational certificates, and alternative strategies. The present paper
keeps only the proof-relevant chain.
Following Connes' reduction of the Riemann Hypothesis to an even-dominance
condition for the Weil quadratic form, we make two independent contributions.
First, for the finite range $100\leq\lambda\leq 1{,}300{,}000$, the Landscape
computation reports negative interval-enclosed Galerkin gaps at 33 grid points.
A 15 May 2026 audit shows that the current sine-block lower-bound construction
still contains a heuristic geometric tail extrapolation; the finite-range data
are therefore treated here as computer-assisted numerical evidence, not as a
theorem-level CAP proof. Second, for the asymptotic regime $\lambda\to\infty$,
we establish via
a common Rayleigh-vector argument that
$\lambda_{\min}(W^+_{\mathrm{prime}}-W^-_{\mathrm{prime}})=-\Theta(\sqrt\lambda)$
on the relevant three-mode Galerkin block. Closing the gap between this
variational statement and the eigenvalue ordering
$\lambda_1^+(\lambda)<\lambda_1^-(\lambda)$ for asymptotic $\lambda$ requires a
separate quantitative lower bound for $\lambda_1^-$; we formulate this as the
\emph{Asymptotic Variational Gap Conjecture} and outline a degenerate
perturbation-theory route in Section~\ref{sec:asymptotic-lower-bound}.
The proof of RH presented here is therefore conditional on (i) the
Connes program (Connes--van Suijlekom real-zero criterion
\cite{ConnesvS2025}, which is fully proven at fixed $L$, together with
the Hurwitz limit-passage of \cite[Sec.~6.4--6.6]{Connes2026}, whose
Fact~6.4 convergence is presented there as a Slepian--Pollak consequence),
and (ii) the
Asymptotic Variational Gap Conjecture for the odd-sector minimum
eigenvalue. The Shift Parity Lemma is unconditional. The
interval-arithmetic certificates as currently implemented use a
heuristic geometric extrapolation for the sin-block lower bound; the
14 May 2026 v1.4 revision documents this honestly in
\cref{sec:k0-strengthened}.

\medskip\noindent
\textbf{v1.5 reformulation (14 May 2026, evening).}  Earlier drafts
expressed the Weil quadratic form on $L^2([-L,L])$ with a real
cosine/sine basis and identified the even/odd parity sectors with
$t\mapsto -t$.  In this formulation, the cosine and sine sectors
share identical symmetric-shift auto-overlaps at every frequency
$k\pi/L$, producing
$\sigma(A_\lambda^+) = \sigma(A_\lambda^-) \cup \{\mu^+_0\}$ exactly
and thus contradicting even dominance.  A comparison against
Connes--van Suijlekom's proven form \cite{ConnesvS2025} shows that
the correct setting is $L^2([0,L])$ with the complex exponential
basis $U_n=L^{-1/2}\exp(2\pi inx/L)$ and parity defined via
$U_n\leftrightarrow U_{-n}$.  In this setting the matrix of the
Weil quadratic form has cosine- and sine-sector diagonals
$D_n\mp S_n$, where $S_n=(\pi n)^{-1}\int_0^L\sin(2\pi ny/L)\,dD(y)$
captures the genuine asymmetry between the two parity sectors.
Numerically $S_n>0$ for the Weil distribution (e.g.\
$S_1\approx 8.63$ at $\lambda=1000$, giving cosine diagonal
$\approx 3.26$ versus sine diagonal $\approx 20.52$), so even
dominance is structurally present in the corrected form.  Section~2
has been rewritten accordingly; downstream sections retain their
legacy structure and require parallel adaptation.

The excluded gauge, total-positivity, and universal-commutator routes
are not used and are referred to the Landscape series.
\end{abstract}

\noindent\textbf{Keywords:} Riemann Hypothesis, Weil quadratic form, even dominance, Connes program, Shift Parity Lemma, computer-assisted evidence.

\newpage
\section{Purpose and Relation to the Landscape Series}

This paper is not a fourth installment of the Landscape series. It is a
proof-only extraction. The Landscape papers~\cite{GeigerLandscapeI,GeigerLandscapeII,GeigerLandscapeIII}
retain the full process: the original gauge-theoretic approach, the obstruction
analysis, non-existence theorems for two attractive but false structural routes,
review reports, and computational handoffs. Those components are useful for
audit and scientific history, but they are not needed in the shortest proof
chain.

The argument presented here uses the following two-part chain:
\[
\begin{gathered}
\text{Shift parity}+\text{IA diagnostics for }100\leq\lambda\leq 1{,}300{,}000\\
\Longrightarrow\ \text{finite-range evidence; theorem-level only under \Cref{thm:even-finite}},
\end{gathered}
\]
together with
\[
\begin{gathered}
\text{frontier dominance}+\text{Asymptotic Variational Gap Conjecture}\\
\Longrightarrow\ \text{Even Dominance for }\lambda\geq\Lambda^\dagger\\
\Longrightarrow\ \text{RH via Connes and Hurwitz (\Cref{thm:rh})}.
\end{gathered}
\]
The first line is diagnostic: the Shift Parity Lemma is unconditional, while
the current finite certificates still require a rigorous odd-sector lower-bound
closure before they become theorem-level CAP. The second line introduces two
conditional inputs: the Connes program (external) and the Asymptotic Variational
Gap Conjecture (\Cref{conj:avg}, internal to this approach;
\Cref{sec:asymptotic-lower-bound}).

\section{Connes' Reduction (v1.5 Reformulation)}

\paragraph{v1.5 reformulation note.}  Earlier drafts of this paper
(up to and including v1.4) presented the Weil quadratic form on
$L^2([-L,L])$ with a real cosine/sine basis and identified the
even/odd sectors with parity under $t\mapsto -t$.  An audit of the
operator definition in May 2026 against Connes--van Suijlekom
\cite{ConnesvS2025} showed that the originally proven form acts on
$L^2([0,L])$ with a complex exponential basis, and that the
relevant parity sectors are defined via the frequency reflection
$U_n\leftrightarrow U_{-n}$ on this space.  Under the simplified
symmetric-shift formulation of earlier drafts the cosine and sine
sectors would have identical spectra (modulo the constant mode),
contradicting the certified gaps; the Connes--van Suijlekom form
admits a genuine asymmetry between cosine and sine sectors,
captured by the quantity $S_n$ introduced below, and it is in this
form that even dominance is mathematically meaningful.  The
present section reformulates the setup to match the proven form;
\Cref{sec:asymptotic-lower-bound} keeps a record of the previous
formulation and the diagnostic chain that led to this revision.

\paragraph{Hilbert space and basis.}  Let $L=\log\lambda$.  Following
\cite[Section~4]{ConnesvS2025}, take the Hilbert space
\[
  \mathcal{H} \;=\; L^2([0,L],\,dx),
\]
with orthonormal Fourier basis
\[
  U_n(x) \;:=\; \frac{1}{\sqrt L}\,\exp\!\Big(\tfrac{2\pi i n x}{L}\Big),
  \qquad x\in[0,L],\ n\in\mathbb{Z},
\]
extended by $0$ outside $[0,L]$.  The involution and convolution are
\[
  f^*(x) \;:=\; \overline{f(-x)},
  \qquad
  (f * g)(y) \;:=\; \int_{\mathbb{R}} f(x)\,g(y-x)\,dx.
\]
For $U_n$ as above, $U_n^*$ is supported on $[-L,0]$ and
$U_n^*(x)=L^{-1/2}\exp(2\pi i n x/L)$ on that interval.

\paragraph{Weil distribution.}  Let $D$ be the Weil distribution on
$[0,L]$,
\begin{equation}
\label{eq:weil-distribution}
  D \;=\; (\log 4\pi + \gamma)\,\delta_0
     \;+\; \frac{e^{y/2}}{2\sinh y}\,\chi_{(0,L]}(y)\,dy
     \;-\; \sum_{p,\,m\geq 1}
           \frac{\log p}{p^{m/2}}\,\delta_{m\log p},
\end{equation}
where $\delta_a$ is the Dirac mass at $a$ and the prime sum runs over
those pairs $(p,m)$ with $m\log p\in[0,L]$, i.e.\ $p^m\le\lambda$.  The
regularised archimedean integrand of \cite[Section~4]{ConnesvS2025}
absorbs the singularity of $e^{y/2}/(2\sinh y)$ at $y=0$ into the
diagonal modification described below.

\paragraph{Quadratic form.}  The Weil quadratic form $Q=Q_\lambda$ on
trigonometric polynomials is
\begin{equation}
\label{eq:quad-form-connesvs}
  \langle f \,|\, g\rangle_Q
  \;:=\; \int_0^L \bigl[(f^* * g)(y) + (f^* * g)(-y)\bigr]\,dD(y).
\end{equation}
By the identity
$(f^**g)(\pm y) = \langle f, T_{\mp y}\,g\rangle$
(with $T_u f(t):=f(t-u)$), this equals
$\int_0^L \langle f, (T_y + T_{-y})g\rangle\,dD(y)$, i.e.\ a
symmetric-shift form integrated against $D$.  We let $A_\lambda$ denote
the lower-bounded self-adjoint operator on $\mathcal{H}$ associated
with $Q$ via the Friedrichs extension.

\paragraph{Matrix elements.}  A direct computation
(\cite[Equation~7]{ConnesvS2025} and Appendix below) gives, for $y\in[0,L]$,
\begin{equation}
\label{eq:matrix-elements-Q}
  (U_m^* * U_n)(y) + (U_m^* * U_n)(-y)
  \;=\;
  \begin{cases}
    2(1-y/L)\cos(2\pi n y/L), & m=n,\\[2pt]
    \displaystyle\frac{\sin(2\pi m y/L) - \sin(2\pi n y/L)}{\pi(n-m)}, & m\neq n.
  \end{cases}
\end{equation}
The matrix of $Q$ in the $\{U_n\}$ basis is therefore
\begin{equation}
\label{eq:matrix-of-Q}
  \langle U_m \,|\, U_n\rangle_Q
  \;=\;
  \int_0^L
  \bigl[(U_m^* * U_n)(y) + (U_m^* * U_n)(-y)\bigr]\,dD(y).
\end{equation}

\paragraph{Diagonal and pair-symmetric off-diagonals.}  Introduce
\begin{equation}
\label{eq:Dn-Sn}
  D_n \;:=\; \int_0^L 2(1 - y/L)\cos\!\Big(\tfrac{2\pi n y}{L}\Big)\,dD(y),
  \qquad
  S_n \;:=\; \frac{1}{\pi n}\int_0^L \sin\!\Big(\tfrac{2\pi n y}{L}\Big)\,dD(y),
\end{equation}
for $n\neq 0$, and $D_0:=\int_0^L 2(1-y/L)\,dD(y)$, $S_0:=0$.  Then
\begin{equation}
\label{eq:U-diag-and-pair}
  \langle U_n\,|\,U_n\rangle_Q \;=\; D_n,
  \qquad
  \langle U_n\,|\,U_{-n}\rangle_Q
  \;=\; \langle U_{-n}\,|\,U_n\rangle_Q
  \;=\; -S_n.
\end{equation}
In particular $D_{-n}=D_n$, so the pair $\{U_n,U_{-n}\}$ acts as an
invariant subspace under the matrix \eqref{eq:matrix-of-Q} only in the
sense that its $2\times 2$ restriction is symmetric in
$\{U_n,U_{-n}\}$; full off-diagonal entries
$\langle U_m|U_n\rangle_Q$ with $m\neq\pm n$ are in general
nonzero but are not used in the diagonal analysis below.

\paragraph{Cosine--sine sector decomposition.}  Define the real
orthonormal basis
\[
  \cos_n \;:=\; \tfrac{1}{\sqrt 2}(U_n + U_{-n}),
  \qquad
  \sin_n \;:=\; \tfrac{1}{i\sqrt 2}(U_n - U_{-n})
  \qquad (n\geq 1),
\]
together with $U_0$ for $n=0$.  Using \eqref{eq:U-diag-and-pair},
\begin{equation}
\label{eq:cos-sin-diagonals}
  \langle \cos_n\,|\,\cos_n\rangle_Q \;=\; D_n - S_n,
  \qquad
  \langle \sin_n\,|\,\sin_n\rangle_Q \;=\; D_n + S_n,
\end{equation}
so the two parity sectors share the same diagonal coefficient $D_n$
but differ by $\pm S_n$.  The asymmetry $2S_n$ is the central object
of this paper: its sign and magnitude determine whether the cosine
sector or the sine sector achieves the smaller diagonal entry, and
hence (modulo controlled off-diagonal corrections) which sector
realises the operator's minimum eigenvalue.

\begin{definition}[Even dominance, Connes--van Suijlekom form]
\label{def:even-dominance}
Let $A_\lambda^{\cos}$ and $A_\lambda^{\sin}$ denote the
$Q_\lambda$-induced operators on the cosine and sine sectors
$\overline{\operatorname{span}}\{U_n+U_{-n}:n\geq 0\}$ and
$\overline{\operatorname{span}}\{U_n-U_{-n}:n\geq 1\}$ respectively.
We say that $\QW_\lambda$ has \emph{even dominance} if
\[
  \lambda_1^{\cos}(\lambda) \;<\; \lambda_1^{\sin}(\lambda)
  \quad\text{with simple, isolated cosine minimum},
\]
where $\lambda_1^{\bullet}(\lambda):=\min\sigma(A_\lambda^{\bullet})$.
We retain the earlier notation $\lambda_1^\pm=\lambda_1^{\cos/\sin}$
for compatibility with later sections.
\end{definition}

\begin{remark}[Numerical asymmetry of $S_n$]
\label{rem:Sn-numerics}
A direct computation of $D_n$ and $S_n$ from
\eqref{eq:Dn-Sn} for the Weil distribution
\eqref{eq:weil-distribution} confirms that $S_n$ is non-negligible:
for $\lambda=1000$ and $n=1$ one finds $D_1\approx 11.89$ and
$S_1\approx 8.63$, so the cosine-sector diagonal is $D_1-S_1\approx 3.26$
while the sine-sector diagonal is $D_1+S_1\approx 20.52$
(script \texttt{\_proof-notes/verify\_connes\_vs\_form.py}).  The
asymmetry $2S_1\approx 17.27$ is of the same order as the diagonal
itself; this is the mechanism by which even dominance arises in the
Connes--van Suijlekom framework and was missed by the legacy
$L^2([-L,L])$ formulation, where the symmetric-shift auto-overlap
forces equal diagonals in the two sectors.
\end{remark}

\paragraph{Downstream migration (v1.5).}  The above
reformulation supersedes the symmetric-shift $L^2([-L,L])$ form
of v1.4.  Subsequent sections (\Cref{sec:parity-matrix} onward) were
written in the legacy framework and require parallel updates: the
$3\times 3$ parity matrix $D_3(r)$ of \Cref{sec:parity-matrix} now
corresponds to the cosine/sine asymmetry at a single shift, with
$D_3(r)$ entries related to $S_n$ contributions from a single prime;
the finite certificates section requires rebuilt CAP matrices in the
$\{U_n\}$ basis; the direct frontier argument carries over once $C_f$
is re-derived as a lower bound on $S_n$ at the frontier-prime
contributions.  The Shift Parity Lemma (\Cref{thm:shift-parity}) and
the obstruction theorems NE-A, NE-B of the Landscape series remain
unaffected.  This migration is partial in the present revision: the
new Section~2 above is the corrected operator definition, while the
downstream propagation is recorded as an open task in
\Cref{sec:k0-strengthened}.  Consequently, the legacy sections below are
retained as diagnostics and route documentation, not as a completed proof in
the corrected $U_n$ framework.

\begin{theorem}[Connes--van Suijlekom real-zero criterion]
\label{thm:connes}
Let a real even distribution define a lower-bounded self-adjoint quadratic-form
operator whose lowest spectral value is simple, isolated, and has an even
eigenfunction $\eta$. Then the Fourier transform $\widehat\eta$ has only real
zeros.
\end{theorem}

\begin{proof}
This is the quadratic-form real-zero criterion proved by Connes and van
Suijlekom~\cite{ConnesvS2025}; Connes' survey applies it to the Weil quadratic
form and identifies even dominance as the remaining spectral condition
\cite{Connes2026}.
\end{proof}

\begin{corollary}[Sufficiency of a divergent sequence]
\label{cor:sequence-suffices}
If even dominance holds for a sequence $\lambda_n\to\infty$, then the Riemann
Hypothesis follows.
\end{corollary}

\begin{proof}
For each $\lambda_n$, \Cref{thm:connes} gives real zeros for the corresponding
truncated Fourier transform $\widehat\eta_{\lambda_n}$. The Hurwitz bridge of
the Connes program~\cite[Fact~6.4 and Sections~6.4--6.6]{Connes2026} establishes:
\begin{enumerate}[label=(\alph*),leftmargin=*]
  \item The truncated kernels $k_{\lambda_n}$ converge locally uniformly to the
  completed Riemann $\Xi$-function on every compact subset of $\mathbb{C}$, at
  rate $O(\lambda_n^{-1/2-\alpha})$ for some $\alpha>0$.
  \item All $k_{\lambda_n}$ and $\Xi$ are entire functions of order one.
  \item By Hurwitz' theorem on zeros of locally uniform limits of holomorphic
  functions, every zero of $\Xi$ in any open set $U\subset\mathbb{C}$ is a limit
  of zeros of $k_{\lambda_n}$ inside $U$ for $n$ large.
\end{enumerate}
Since each $k_{\lambda_n}$ has only real zeros (\Cref{thm:connes}), the zeros of
$\Xi$ are real as well. The non-trivial zeros of $\zeta(s)$ correspond bijectively
to zeros of $\Xi$, hence they lie on the critical line~$\Re(s)=1/2$.
\end{proof}

\begin{remark}[Analytical structure of the Hurwitz bridge]
\label{rem:hurwitz-analytical}
For transparency, we spell out the analytical chain underlying step~(a) above.
The truncated kernel is
\[
  k_\lambda(z)
  = \int_{-L}^{L} K_\lambda(t)\,e^{izt}\,dt, \qquad L = \log\lambda,
\]
where $K_\lambda$ is the even Weil spectral kernel restricted to $[-L,L]$.
As a finite Fourier integral over a bounded interval, $k_\lambda$ belongs to
the Paley--Wiener class $\mathrm{PW}_L$: it extends to an entire function
satisfying $|k_\lambda(z)|\leq C_\lambda\,e^{L|\Im z|}$.

A self-contained proof of locally uniform convergence $k_\lambda\to\Xi$
would require three ingredients:
\begin{enumerate}[label=(\roman*),leftmargin=*]
  \item \emph{$L^2$-rate control}: $\|K_\lambda - K\|_{L^2(\mathbb{R})}
    = O(\lambda^{-1/2-\alpha})$ for some $\alpha>0$,
    where $K$ is the full (untruncated) Weil kernel.
    This requires Schwartz-class regularity of $K_\lambda - K$
    at the stated rate --- the non-trivial analytical input.
  \item \emph{Locally uniform convergence from $L^2$ and Paley--Wiener}:
    on any compact $\mathcal{K}\subset\mathbb{C}$, PW membership of $k_\lambda$
    gives $|k_\lambda(z)|\leq e^{LR}\|K_\lambda\|_{L^2}$
    (with $R = \max_{z\in\mathcal{K}}|\Im z|$), so the family is
    locally bounded; normality (Montel's theorem) and $L^2$-convergence
    then yield locally uniform convergence on $\mathcal{K}$.
  \item \emph{Order one}: both $k_\lambda$ and $\Xi$ have order one,
    as required for Hurwitz' theorem to transfer all zeros.
\end{enumerate}
Connes establishes step~(i) via the spectral theory of his Hilbert space in
\cite[Fact~6.4 and Sections~6.4--6.6]{Connes2026};
steps~(ii) and~(iii) are classical.
The present paper uses Connes' result for step~(i) directly and does not
provide an independent derivation.
\end{remark}

\begin{remark}[Conditioning and hypotheses]
\label{rem:conditioning}
The present proof relies on the quadratic-form real-zero criterion of
Connes and van Suijlekom~\cite{ConnesvS2025} and on Connes' identification
of even dominance as the remaining spectral condition~\cite{Connes2026}.
Both results are currently available as arXiv preprints. Our argument is
therefore conditional on those results until they are independently verified
or published in a peer-reviewed journal.
The hypotheses of \Cref{thm:connes} require: (i)~that $A_\lambda$ be bounded
below, which holds because the archimedean kernel is integrable and the
prime sum converges absolutely; (ii)~that the lowest spectral value be simple
and isolated --- this is the content of the even-dominance condition proved
below; and (iii)~locally uniform convergence of the Fourier transforms of the
truncated eigenfunctions, established in~\cite[Section~4]{Connes2026}.
\end{remark}

\section{The Parity Matrix of a Single Shift}
\label{sec:parity-matrix}

\paragraph{v1.5 reformulation note.}  This section, originally written
in the legacy $L^2([-L,L])$ cosine/sine framework, has been retained
in its legacy form below (Subsection~\ref{subsec:parity-legacy}) and
augmented with the Connes--van Suijlekom analog
(Subsection~\ref{subsec:parity-cvs}).  In the legacy framework the
matrix $D_3(r)$ encodes the cos--sin parity difference at a single
shift directly; in the Connes--van Suijlekom framework the analog is
the single-shift contribution to the asymmetry coefficient $S_n$
introduced in \eqref{eq:Dn-Sn}.  Both objects encode the same
content: a single prime shift, after symmetrisation in the Weil
form, contributes \emph{differently} to the cosine and sine
sectors, with the difference quantified by a single trigonometric
expression in the normalised shift $r=u/L$.  The legacy formulation
is preserved for continuity of notation in
Sections~\ref{sec:finite-certificates}--\ref{sec:tail}, which still
read in legacy form; the migration of those sections is recorded as
an open task in \cref{sec:k0-strengthened}.

\subsection{Connes--van Suijlekom single-shift asymmetry}
\label{subsec:parity-cvs}

In the corrected setup of Section~2, a single Dirac mass at $y\in(0,L)$
with weight $w$ in the Weil distribution $D$ contributes to the
diagonal matrix elements as
\[
  \Delta D_n \;=\; w\cdot 2(1-y/L)\cos(2\pi n y/L),
  \qquad
  \Delta S_n \;=\; \frac{w}{\pi n}\sin(2\pi n y/L)
  \quad (n\geq 1).
\]
The Cos- and Sin-sector diagonal contributions are therefore
$\Delta D_n - \Delta S_n$ and $\Delta D_n + \Delta S_n$ respectively
(cf.~\eqref{eq:cos-sin-diagonals}), and the parity asymmetry from this
single shift at mode $n$ is
\begin{equation}
\label{eq:single-shift-asymmetry}
  -2\Delta S_n \;=\; -\frac{2w}{\pi n}\sin(2\pi n y/L)
  \quad
  (\text{cosine-minus-sine diagonal contribution}).
\end{equation}

\begin{lemma}[Single-shift asymmetry; v1.5 analogue of Shift Parity]
\label{lem:shift-parity-cvs}
Let $r:=y/L\in(0,1]$ be the normalised shift.  For a single Dirac mass
$\delta_y$ with the Weil prime-sign weight $w=-(\log p)/p^{m/2}$
(i.e.\ $w<0$), the single-shift cosine-minus-sine asymmetry at mode
$n\geq 1$ is
\[
  -2\Delta S_n \;=\; \frac{2\log p}{\pi n\cdot p^{m/2}}\,
  \sin(2\pi n r).
\]
This is \emph{strictly negative} (i.e.\ favours the cosine sector) iff
$\sin(2\pi n r)<0$, which holds on the union of intervals
\[
  r \;\in\; \bigcup_{k=0}^{n-1}\Bigl(\tfrac{k}{n}+\tfrac{1}{2n},\,
                                       \tfrac{k+1}{n}\Bigr).
\]
In particular, for $n=1$ the favourable range is
$r\in(\tfrac{1}{2},1)$, i.e.\ $\sqrt\lambda<p^m\le\lambda$; this is
the analogue of the legacy ``Frontier'' regime, except that the
favourable range is the half-frontier $r>1/2$ rather than the
neighbourhood of $r=1$.
\end{lemma}

\begin{proof}
Direct substitution of $w=-(\log p)/p^{m/2}$ into
\eqref{eq:single-shift-asymmetry} gives the displayed formula.  The
sign condition $-2w\sin(2\pi n r)<0$ becomes $\sin(2\pi nr)<0$ since
$w<0$, whose solution set on $r\in(0,1]$ is the union of intervals
displayed.  The frontier identification follows from
$y=m\log p$ and $r=y/L=m\log p/\log\lambda$.
\end{proof}

\begin{remark}[Multi-mode coverage]
\label{rem:multi-mode-coverage}
For any $r\in(0,1)$ that is irrational with respect to $\mathbb{Q}$,
the union $\bigcup_n\{n:\sin(2\pi n r)<0\}$ is infinite; in
particular, for every $r$ outside a measure-zero set, the single
shift contributes favourably to infinitely many modes.  The
aggregate over modes is not simply additive in eigenvalue space, but
the Cesaro-style averaging of $-2\Delta S_n$ over $n$ matches the
classical Walfisz-type Mertens estimates entering the asymptotic
analysis of Sections~\ref{sec:frontier}--\ref{sec:tail} (in their
v1.5 reformulation).
\end{remark}

\begin{lemma}[$L$-independence of single-shift asymmetry]
\label{lem:L-independence-cvs}
The quantity $-2\Delta S_n$ depends on the normalised shift $r=y/L$
and the mode $n$, but not on $L$ separately, up to the overall
prime-weight factor $\log p/p^{m/2}$ which is determined by the
shift location $y=m\log p$ and the choice of $L=\log\lambda$.
\end{lemma}

\begin{proof}
Immediate from \eqref{eq:single-shift-asymmetry}: the factor
$\sin(2\pi n r)$ depends only on $r$.
\end{proof}

\subsection{Legacy parity matrix (retained for continuity)}
\label{subsec:parity-legacy}

\paragraph{Note.}  The remainder of this section is the original
$L^2([-L,L])$ Cos/Sin-basis formulation that supplied the technical
content of Sections~\ref{sec:finite-certificates}--\ref{sec:tail}.
It is retained here in legacy form for backward reference;
see~\Cref{subsec:parity-cvs} for the Connes--van Suijlekom analog.
Translating downstream sections to the Connes--van Suijlekom basis
is recorded as an open task in \cref{sec:k0-strengthened}.

The legacy construction begins with a finite algebraic fact. Let
\[
  \varphi_0(t)=\frac{1}{\sqrt{2L}},
  \qquad
  \varphi_n(t)=\frac{\cos(n\pi t/L)}{\sqrt L}\;(n\geq 1),
  \qquad
  \psi_n(t)=\frac{\sin((n+1)\pi t/L)}{\sqrt L}
\]
denote the $L^2([-L,L])$-normalized even and odd basis functions. For a normalized shift $r=s/L$ define
the $N\times N$ difference matrix
\[
  D_{nm}(r)=
  \langle\varphi_n,(S_{rL}+S_{-rL})\varphi_m\rangle
  -
  \langle\psi_n,(S_{rL}+S_{-rL})\psi_m\rangle .
\]

\begin{lemma}[$L$-independence (legacy)]
\label{lem:L-independence}
$D_N(r)$ depends only on $r$, not on $L$ separately.
\end{lemma}

\begin{proof}
Substitute $x=t/L$ in each overlap integral. The normalization factors cancel
the scale factor from $dt=L\,dx$, leaving integrals over fixed intervals with
argument $r$ only.
\end{proof}

\begin{proposition}[Closed diagonal entries]
\label{prop:diagonal}
For $L=1$ the diagonal entries satisfy, for $n\geq 1$:
\[
  D_{nn}(r)
  =(2-r)\bigl[\cos(n\pi r)-\cos((n+1)\pi r)\bigr]
  -\frac{\sin(n\pi r)}{n\pi}
  -\frac{\sin((n+1)\pi r)}{(n+1)\pi}.
\]
For $n=0$, a direct overlap-length computation gives
\[
  D_{00}(r)=(2-r)(1-\cos\pi r)-\frac{\sin\pi r}{\pi}.
\]
\end{proposition}

\begin{theorem}[Shift Parity Lemma]
\label{thm:shift-parity}
For every $N\geq 3$ and every $r\in(0,2)$,
\[
  \lambda_{\min}(D_N(r))<0.
\]
\end{theorem}

\begin{proof}
By the Cauchy interlacing theorem, $\lambda_{\min}(D_{N+1})\leq\lambda_{\min}(D_N)$,
so it suffices to prove the $N=3$ case. The trace telescopes via
\Cref{prop:diagonal}:
\[
  \Tr(D_3(r))=(2-r)(1-\cos 3\pi r)
  -\frac{2\sin\pi r}{\pi}
  -\frac{\sin 2\pi r}{\pi}
  -\frac{\sin 3\pi r}{3\pi}.
\]
The determinant $\det(D_3(r))$ is an explicit trigonometric expression in~$r$,
computed from the closed-form entries of \Cref{prop:diagonal}.
We use a two-case argument.

\smallskip\noindent\textbf{Case~1} ($\det D_3(r)<0$, approximately $93.3\%$ of
$(0,2)$): the eigenvalues have mixed signs, so $\lambda_{\min}<0$.

\smallskip\noindent\textbf{Case~2} ($\det D_3(r)\geq 0$, the interval
$R_2\approx[0.597,0.729]$): in this region, $\Tr(D_3(r))<0$
($\max_{R_2}\Tr\approx -0.007$). Since the trace equals the sum of eigenvalues,
not all three can be non-negative, hence $\lambda_{\min}<0$.

\smallskip\noindent\textbf{Completeness.} The determinant has exactly two zeros
in $(0,2)$, enclosed by interval arithmetic (\texttt{mpmath}, $50$-digit
precision) as
\[
  r_1^{\det}\in[0.59651,0.59653],\qquad r_2^{\det}\in[0.72928,0.72930].
\]
The trace zeros are similarly enclosed:
\[
  r_2^{\Tr}\in[0.58760,0.58762],\qquad r_3^{\Tr}\in[0.73023,0.73025].
\]
The interleaving
\[
  r_2^{\Tr}<r_1^{\det}<r_2^{\det}<r_3^{\Tr}
\]
holds rigorously (interval separation $>0.008$ at each adjacent pair),
ensuring $R_2\subset\{\Tr<0\}$. Cases~1 and~2 cover $(0,2)$ without gap.
See~\cite{GeigerLandscapeII} for the verification scripts.
\end{proof}

\begin{remark}
The lemma is independent of the zeta function. Number theory enters only after
the matrices $D_3(m\log p/\log\lambda)$ are weighted by
$(\log p)p^{-m/2}$ and summed over primes.
\end{remark}

\section{Finite Certificates}
\label{sec:finite-certificates}

\paragraph{v1.5 status of the CAP construction.}  The 33 certificates
described below were computed in the legacy $L^2([-L,L])$ cosine/sine
Galerkin basis with symmetric-shift operators.  Within that basis,
both the cosine-sector upper bound on $\lambda_1^+$ (via interval
arithmetic) and the sine-sector ``lower bound'' on $\lambda_1^-$ (via
heuristic geometric extrapolation of column norms) were treated as
rigorous and produced a strictly negative \texttt{cert\_gap} at every
grid point.  The v1.5 audit (Section~\ref{sec:k0-cvs-resolution};
companion notes \texttt{K0\_CAP\_AUDIT.md} and
\texttt{K0\_CAP\_CONNESVS\_RESULTS.md}) shows two structural issues
with this construction:
\begin{enumerate}[label=(\arabic*),leftmargin=*]
  \item The cosine--sine equality of symmetric-shift auto-overlaps at
    each frequency makes $\sigma(A^+) = \sigma(A^-) \cup \{\mu^+_0\}$
    exactly in the legacy operator form, contradicting the certified
    gaps.  This is resolved in the corrected Connes--van Suijlekom
    form of Section~2 (in which the cosine and sine sectors of the
    operator have genuinely distinct diagonals $D_n\mp S_n$ via the
    frequency-reflection $U_n\leftrightarrow U_{-n}$), but it
    invalidates the literal interpretation of the certificates as
    proving the original $\lambda_1^+ < \lambda_1^-$ on $L^2([-L,L])$.
  \item The sine-block ``lower bound'' is not mathematically
    rigorous: it is computed as the smallest Galerkin eigenvalue minus
    a heuristic geometric extrapolation of off-block column norms.
    A rigorous lower bound would require Lehmann--Maehly,
    Kato--Temple with an a priori lower bound on $\lambda_2$, or a
    quantitative Schur--complement argument; an analysis of the
    obstructions is given in \texttt{K0\_LEHMANN\_ANALYSIS.md}.
\end{enumerate}
The certificates therefore stand as \emph{computer-assisted numerical
evidence} for even dominance, consistent with the corrected
Connes--van Suijlekom structural picture, but not as a fully rigorous
CAP-style proof in the formal sense of computer-assisted proofs in
analysis.  A v1.5 reformulation of the CAP in the $\{U_n\}$ basis with
a rigorous lower-bound technique is recorded as an open task.

\bigskip

Let
\[
  \Delta(\lambda)=\lambda_1^+(\lambda)-\lambda_1^-(\lambda).
\]
The Landscape computation supplies interval-arithmetic certificates for
$33$ values of $\lambda$ in the interval
\[
  100\leq\lambda\leq 1{,}300{,}000.
\]
\begin{sloppypar}
Certificates are computed using Python's \texttt{mpmath} library in interval-arithmetic
mode (\texttt{mpmath.iv}, 50-digit precision). The even sector uses a $4\times 4$
Galerkin matrix ($N_{\cos}=4$); the odd sector uses $N_{\sin}=N_0\in[40,90]$
(see \Cref{rem:galerkin-asymmetry} below). Matrix entries are evaluated as rigorous
intervals: prime shift integrals have width $<10^{-12}$, and the archimedean kernel
integral is enclosed via composite Gauss--Legendre quadrature (width~$<10^{-3}$).
If the sine-block lower-bound construction is replaced by a rigorous bound,
a negative reported gap would imply even dominance; at present it is recorded
as diagnostic numerical evidence.
\end{sloppypar}

\begin{remark}[Galerkin-asymmetry justification]
\label{rem:galerkin-asymmetry}
The asymmetric truncation $N_{\cos}=4$ vs.\ $N_{\sin}\geq 40$ is justified by
the structural difference between the two sectors observed in
\cite[Section~2.4 and Remark on low-mode dominance]{GeigerLandscapeII}: the even
block converges rapidly because the leading even mode $n=1$ has Galerkin diagonal
$(W_{\mathrm{prime}})_{11}^+ \sim -\sqrt\lambda$ (driven by the auto-overlap value
$h_1(1)=-1/4$ at the Laplace endpoint, see \Cref{def:auto-overlap}), with the
$k\times k$ minimum eigenvalue stabilizing at $k=4$ to within $10^{-2}$ for every
tested $\lambda$. The odd sector has no analogous structural amplification, so
$N_0=40$--$90$ modes are required for an equivalent enclosure. Crucially, the
Landscape certifier reports
\(
  \lambda_{\min}(\mathrm{QW}_{\sin}^{(N)})\;>\;\lambda_{\min}(\mathrm{QW}_{\cos}^{(4)})
\)
for \emph{all} $N\leq 90$ at every tested $\lambda$, supporting, but not yet
proving, that the gap is genuine and not a truncation artifact.
\end{remark}

\begin{remark}[Truncation safety factor]
\label{rem:truncation-safety}
The reported Galerkin/tail truncation error of the Landscape certifier at $\lambda=100$
is at most $0.061$ (Frobenius interval radius $\leq 0.0025$ in the even block, total
odd-sector tail $\leq 0.059$ via Cauchy interlacing plus a float64 rounding margin),
while $|\mathrm{cert\_gap}|=2.25$ -- a safety factor $\geq 34\times$. For $\lambda\geq 200$
the diagnostic safety factor exceeds $65\times$ and grows monotonically with $\lambda$
throughout the sampled range~\cite[Section~2.9 and certificate metadata]{GeigerLandscapeII}.
\end{remark}

\begin{proposition}[Finite certificate diagnostic]
\label{prop:finite-certificates}
For every certified grid point $\lambda_k$ in the Landscape table,
\[
  \mathrm{cert\_gap}(\lambda_k)<0.
\]
Moreover, the reported even-sector spectral gap is positive at every grid
point. These data are numerical evidence; they become a rigorous finite bridge
only after the sine-block lower-bound construction is validated or replaced in
the corrected Connes--van Suijlekom framework.
\end{proposition}

\begin{proof}
This is a computer-assisted diagnostic output from the Landscape project
\cite{GeigerLandscapeII}. The relevant certificate files are stored in the Landscape
repository under \texttt{results/certificates/}. Representative values are:
\[
\begin{array}{r|r|r|r}
\lambda & \#\{p\leq\lambda\} & \lambda_1^+ \text{ upper bound} &
\mathrm{cert\_gap}\\\hline
100 & 25 & -11.07 & -2.25\\
500 & 95 & -34.13 & -7.72\\
10{,}000 & 1{,}229 & -216.22 & -42.37\\
320{,}000 & 27{,}608 & -1220.97 & -241.75\\
1{,}300{,}000 & 100{,}021 & -2503.78 & -475.83
\end{array}
\]
All reported intervals have negative gap. The same certificate set gives a
positive intra-even gap, supporting the absence of even-sector level crossing at
the sampled points.
\end{proof}

\begin{remark}[Certificate structure]
\label{rem:cert-structure}
Each diagnostic certificate for a given $\lambda_k$ consists of:
(a)~the $N_{\cos}\times N_{\cos}$ even-sector Galerkin matrix with interval-arithmetic
entries;
(b)~the $N_{\sin}\times N_{\sin}$ odd-sector Galerkin matrix;
(c)~a rigorous eigenvalue enclosure
$[\underline\lambda_1^+,\overline\lambda_1^+]$ for the minimum even eigenvalue;
(d)~a candidate lower bound $\underline\lambda_1^-$ for the minimum odd eigenvalue,
currently relying on the heuristic tail extrapolation described above;
(e)~the reported gap $\overline\lambda_1^+-\underline\lambda_1^-<0$.
For instance, at $\lambda=100$: the $4\times 4$ even matrix has minimum eigenvalue
in $[-11.08,-11.06]$, the odd diagnostic lower-bound output lies above $-8.82$,
giving reported gap $<-2.24$. All certificates, including the interval-arithmetic evaluation scripts,
are archived in~\cite{GeigerLandscapeII}.
\end{remark}

\section{Frontier Dominance}
\label{sec:frontier}

\paragraph{v1.5 reformulation summary.}  In the Connes--van Suijlekom
form of Section~2, the analog of frontier dominance is the asymptotic
positivity of the cos/sin sector asymmetry $S_n$ for each mode~$n$.
For $n=1$, an Abel--summation and Laplace asymptotic computation
gives
\begin{equation}
\label{eq:S1-asymptotic-cvs}
  S_1(\lambda) \;\sim\; \frac{8\sqrt\lambda}{\log\lambda}
  \quad (\lambda\to\infty),
\end{equation}
\begin{sloppypar}
with the more accurate finite-$L$ form
$S_1(\lambda) \approx 8\sqrt\lambda/[L(1+16\pi^2/L^2)]$ where
$L=\log\lambda$.  Numerically verified to within $5{-}10\%$ on
$\lambda\in\{100,300,1000\}$ against the script
\texttt{verify\_connes\_vs\_form.py}, see
\texttt{K0\_FRONTIER\_SN\_ASYMPTOTIC.md} for the
derivation.
\end{sloppypar}  In this v1.5 form the cosine--minus--sine diagonal gap at
mode~$n=1$ is
\[
  \langle\sin_1\,|\,Q\,|\,\sin_1\rangle
    - \langle\cos_1\,|\,Q\,|\,\cos_1\rangle
  \;=\; 2 S_1(\lambda)
  \;\sim\; \frac{16\sqrt\lambda}{\log\lambda},
\]
i.e.\ a logarithmically shrinking fraction of the leading
$\sqrt\lambda$ scale.  This is positive asymptotic evidence for even dominance
on the $O(\sqrt\lambda/\log\lambda)$ scale, but the eigenvalue-ordering
statement still requires the Asymptotic Variational Gap or an equivalent
lower-bound closure.  The Legacy
direct-frontier construction below is retained for backward
compatibility with the rest of the manuscript; the analog $S_1$
construction can be developed in the v1.5 form by replacing
$D_3(r)$ with the single-shift asymmetry function
$f(r):=\sin(2\pi r)/\pi$ (cf.\ \cref{lem:shift-parity-cvs}).

\bigskip

The Legacy asymptotic proof uses the fact that primes near the moving frontier
$p\sim\lambda$ dominate the parity sum.

Define the low-mode cumulative difference matrix
\[
  D_{\mathrm{tot}}(\lambda)=
  \sum_{\substack{p,m\geq1\\0<m\log p/\log\lambda<2}}
  \frac{\log p}{p^{m/2}}\,
  D_3\!\left(\frac{m\log p}{\log\lambda}\right).
\]
Split it into
\[
  D_{\mathrm{tot}}(\lambda)=D_f(\lambda)+D_n(\lambda),
\]
where $D_f$ contains the first-order shifts with
$p\in[\lambda^{0.7},\lambda]$ and $D_n$ contains the remaining terms.

\begin{lemma}[Common frontier Rayleigh vector]
\label{lem:common-vector}
Let $e_1=(0,1,0)^T$. For all $r\in[0.7,1]$,
\[
  e_1^T D_3(r)e_1\leq -C_f,
  \qquad C_f=0.4685.
\]
\end{lemma}

\begin{proof}
By \Cref{prop:diagonal},
\[
  e_1^T D_3(r)e_1
  =(2-r)(\cos\pi r-\cos2\pi r)
  -\frac{\sin\pi r}{\pi}
  -\frac{\sin2\pi r}{2\pi}.
\]
Interval evaluation of this one-variable trigonometric function on $[0.7,1]$
gives maximum $<-0.4685$, attained at the left endpoint. Thus the same vector
$e_1$ is a valid negative Rayleigh vector for the whole frontier window.
\end{proof}

\begin{lemma}[Frontier lower bound]
\label{lem:frontier}
For the frontier block,
\[
  \lambda_{\min}(D_f(\lambda))
  \leq
  -C_f W_f(\lambda),
  \qquad
  W_f(\lambda)=
  \sum_{\lambda^{0.7}\leq p\leq\lambda}\frac{\log p}{\sqrt p}.
\]
Furthermore,
\[
  W_f(\lambda)=2\sqrt\lambda(1-\lambda^{-0.15})+O(1).
\]
\end{lemma}

\begin{proof}
Use the fixed Rayleigh vector $e_1$ from \Cref{lem:common-vector}:
\[
  \lambda_{\min}(D_f)
  \leq e_1^T D_f e_1
  \leq -C_f\sum_{\lambda^{0.7}\leq p\leq\lambda}\frac{\log p}{\sqrt p}.
\]
The asymptotic estimate follows by partial summation~\cite{IwaniecKowalski2004}
from the Prime Number Theorem in the form $\theta(x)\sim x$.
\end{proof}

\begin{lemma}[Non-frontier norm bound]
\label{lem:nonfrontier}
There is an absolute constant $C_n$ such that
\[
  \norm{D_n(\lambda)}_{\mathrm{op}}
  \leq C_n\bigl(2\lambda^{0.35}+O(\log\lambda)\bigr).
\]
One may take $C_n=2.2847$ from the interval-certified bound
$\sup_{0<r<2}\norm{D_3(r)}_{\mathrm{op}}\leq2.2847$.
\end{lemma}

\begin{proof}
The first non-frontier part, $p<\lambda^{0.7}$ with $m=1$, has total weight
\[
  \sum_{p\leq\lambda^{0.7}}\frac{\log p}{\sqrt p}
  =2\lambda^{0.35}+O(1)
\]
by partial summation with the unconditional estimate
$|\theta(x)-x|\leq 0.006788\,x/\log x$ for $x\geq 89{,}967{,}803$~\cite{Dusart2010};
for smaller $x$ the sum is bounded by direct computation.
The higher-shift terms satisfy
\[
  \sum_{p\leq\lambda}\sum_{m\geq2}\frac{\log p}{p^{m/2}}
  =O(\log\lambda),
\]
with the $m=2$ term giving the logarithmic growth and $m\geq3$ bounded by
convergent prime sums. Multiplying by the uniform operator-norm bound for
$D_3(r)$ gives the claim.
\end{proof}

\begin{proposition}[Legacy variational frontier dominance]
\label{prop:frontier-dominance}
There exists an explicit $\Lambda_*$ such that for all $\lambda\geq\Lambda_*$,
\[
  \lambda_{\min}(D_{\mathrm{tot}}(\lambda))<0.
\]
With the constants above, $\Lambda_*$ lies below the certified endpoint
$1{,}300{,}000$ after the explicit error terms are evaluated.
\end{proposition}

\begin{proof}
We argue via the Rayleigh quotient at the fixed common test vector $e_1=(0,1,0)^T$
together with the Cauchy--Schwarz operator-norm bound (no eigenvalue additivity
is invoked). The variational inequality gives
\[
  \lambda_{\min}(D_{\mathrm{tot}})
  \;\leq\;
  e_1^T D_{\mathrm{tot}}\, e_1
  \;=\;
  e_1^T D_f\, e_1 + e_1^T D_n\, e_1.
\]
For the frontier block, \Cref{lem:common-vector,lem:frontier} give
$e_1^T D_f e_1 \leq -C_f W_f(\lambda) = -2C_f\sqrt\lambda(1-\lambda^{-0.15})+O(1)$.
For the non-frontier block, Cauchy--Schwarz yields
$|e_1^T D_n e_1| \leq \|e_1\|^2\,\norm{D_n}_{\mathrm{op}} = \norm{D_n}_{\mathrm{op}}$,
which is bounded by \Cref{lem:nonfrontier}. Combining,
\[
  \lambda_{\min}(D_{\mathrm{tot}})
  \;\leq\;
  -2C_f\sqrt\lambda(1-\lambda^{-0.15})
  + 2C_n\lambda^{0.35}
  + O(\log\lambda).
\]
The negative term is of order $\sqrt\lambda$, the positive non-frontier term of
order $\lambda^{0.35}$, so the right-hand side is strictly negative for all
sufficiently large $\lambda$. Evaluating the displayed bound with the interval-certified
constants $C_f=0.4685$, $C_n=2.2847$ and the explicit PNT/Mertens remainders
gives an effective threshold $\Lambda_*\leq 1{,}300{,}000$; see
\cite[Proposition~A6 (Direct Frontier-Dominance), Remark on direct-proof
ingredients]{GeigerLandscapeII}. The asymptotic regime $\lambda\geq\Lambda_*$ and
the finite diagnostic range $100\leq\lambda\leq 1{,}300{,}000$ of
\Cref{prop:finite-certificates} are numerically compatible but do not yet give
a theorem-level bridge without the lower-bound repair.
\end{proof}

\begin{remark}[Computational improvement of $\Lambda_*$]
\label{rem:lambda-star-improved}
The bound $\Lambda_*\leq 1{,}300{,}000$ above is conservative: it is the right
endpoint of the certified interval. The exact Rayleigh gap
\[
  g(\lambda) = -C_f\,W_f(\lambda) + C_n\,W_n(\lambda)
\]
(with $W_f$, $W_n$ the precise frontier and non-frontier prime-sum weights,
computed by exact prime enumeration) was evaluated as follows.

\noindent\emph{(i) Step-1 exhaustive check $[171{,}329,\,182{,}000]$:}
Every integer $\lambda$ in this range was checked; $g(\lambda)<0$ throughout,
with the first negative value $g(171{,}329)=-0.0052$.

\noindent\emph{(ii) Step-1000 coarse scan $[182{,}000,\,300{,}000]$:}
All sampled values give $g\leq -3.5$.
In a 500-step window at most two prime-exit events can occur, each increasing $g$
by at most $(C_f+C_n)\log p/\sqrt{p}\leq 0.35$; the total upward perturbation
$\leq 0.70$ leaves a margin of $\geq 2.8$ below zero.
No unsampled integer in $[182{,}000, 300{,}000]$ can be positive.

\noindent\emph{(iii) Asymptotic regime $\lambda\geq 300{,}000$:}
$g\leq -43$ throughout (coarse scan, step $10{,}000$), consistent with the
$-\Theta(\sqrt\lambda)$ dominance of the asymptotic formula.

\noindent Combining (i)--(iii) with the finite diagnostic range
$[100,\,1{,}300{,}000]$ gives numerical support for the threshold
$\Lambda_*\leq 175{,}000$ in the variational model.
The verification script \path{verify_lambda_star.py} (float64, safety margins
$\gg 10^{-10}$) is archived in the companion repository at \projectrepo.
\end{remark}

\section{Tail Control and Lifting to the Weil Form}
\label{sec:tail}

\paragraph{v1.5 reformulation note.}  This section, written in the
legacy $L^2([-L,L])$ basis, controls the higher-mode tail
contributions in that framework.  In the Connes--van Suijlekom form
of Section~2, the analogous task is to control the off-diagonal
contributions of the matrix $\langle U_m|Q|U_n\rangle$ for
$|m|,|n|>N$ where $N$ is the Galerkin cutoff, plus the regularised
archimedean integrand.  The estimates below suggest the corresponding
objects after replacing the relevant overlap functions by
$\sin(2\pi m y/L) - \sin(2\pi n y/L)$ in place of $D_3(r)$ and replacing the
Schur-complement bound by one for the off-block norm of the
Connes--van Suijlekom matrix.  They are not used here as migrated
$U_n$-proofs. Migration of the precise inequalities
is open and is recorded with the analytical roadmap in
\texttt{K0\_SN\_SUBLEADING.md}.  In particular, the subleading
asymptotic $S_n \sim 8\sqrt\lambda/[L(1+16n^2\pi^2/L^2)]$ replaces
the leading $\Theta(\sqrt\lambda)$ contribution of the legacy
direct-frontier bound, giving an $O(\sqrt\lambda/\log\lambda)$ scale
for the Cos--Sin diagonal gap.

\bigskip

\Cref{prop:frontier-dominance} establishes the legacy variational statement
$\lambda_{\min}(D_{\mathrm{tot}}(\lambda))<0$ in the three-mode difference
block. The full Weil operator $A_\lambda$ acts on infinitely many modes and
includes the archimedean contribution. The estimates below document why the
legacy low-mode calculation was numerically stable; in v1.5 they are not a
completed proof that the corrected Connes--van Suijlekom $U_n$ operator
preserves an even--odd eigenvalue gap.

\begin{remark}[Archimedean contribution]
\label{rem:archimedean}
The archimedean kernel $e^{u/2}/(2\sinh u)$ in~\eqref{eq:quad-form-connesvs} is integrable
on $(0,\infty)$, and the sector-independent term $-2e^{-u/2}\norm{f}^2$ cancels in the
even--odd difference. The remaining archimedean contribution involves shift overlap
differences $D_N(u/L)$ weighted against the fixed kernel. Since
$\norm{D_3(r)}_{\mathrm{op}}\leq 2.2847$ uniformly and the kernel integral is a
constant independent of $\lambda$, the archimedean even--odd contribution satisfies
$|\Delta_{\mathrm{arch}}|\leq C_{\mathrm{arch}}$ for an absolute constant. As the
prime contribution is $\Theta(\sqrt\lambda)$ by \Cref{lem:frontier}, the archimedean
part is negligible for large $\lambda$.
\end{remark}

\begin{definition}[Auto-overlap profile]
\label{def:auto-overlap}
For the even basis function $\varphi_n$ and normalized shift $u=\delta/L\in[0,2]$,
define $h_n(u):=\tfrac14\langle\varphi_n,(S_{uL}+S_{-uL})\varphi_n\rangle$.
The overlap integral evaluates via the substitution $x=t/L$ and the
product-to-sum identity
$\cos\alpha\cos\beta=\tfrac12[\cos(\alpha-\beta)+\cos(\alpha+\beta)]$:
\begin{align*}
  \langle\varphi_n,(S_{uL}+S_{-uL})\varphi_n\rangle
  &= 2\int_{u-1}^{1}\cos(n\pi x)\cos(n\pi(x-u))\,dx \\
  &= \cos(n\pi u)(2-u)
    + \frac{\sin(n\pi(2-u))-\sin(n\pi(u-2))}{2n\pi} \\
  &= (2-u)\cos(n\pi u) - \frac{\sin(n\pi u)}{n\pi},
\end{align*}
where the last step uses $\sin(n\pi(2-u))=-\sin(n\pi u)$ for integer~$n$.
Dividing by~$4$:
\[
  h_n(u) = \tfrac{1}{2}(1-\tfrac{u}{2})\cos(n\pi u)
           - \frac{\sin(n\pi u)}{4n\pi},
  \qquad n\geq 1,
\]
with $h_0(u)=(1-u/2)/2$ (from $\langle\varphi_0,(S_{uL}+S_{-uL})\varphi_0\rangle
=2-u$ by direct overlap-length computation).
At the Laplace endpoint $u=1$: $h_n(1)=(-1)^n/4$ for all $n\geq 0$.
The normalization factor~$\tfrac14$ is chosen so that the prime-weighted diagonal
satisfies $(W_{\mathrm{prime}})_{nn}=\sum_p(\log p)\,p^{-1/2}\cdot 4h_n(\log p/L)$,
consistent with the Weil form~\eqref{eq:quad-form-connesvs}.
\end{definition}

\begin{lemma}[Leading-mode cancellation]
\label{lem:leading-cancel}
Define the symmetrized off-diagonal overlap profiles
$F_j(u):=S^-_{\mathrm{sym}}(0,j;u)$ (odd sector) and
$G_j(u):=S^+_{\mathrm{sym}}(1,j';u)$ (even sector) for the leading modes.
Their closed forms are:
\begin{align}
  S^+_{\mathrm{sym}}(1,0;u)
    &= \frac{\sqrt{2}\sin(\pi u)}{\pi}, \label{eq:Scos10} \\
  S^+_{\mathrm{sym}}(1,2;u)
    &= \frac{2(4\cos\pi u - 1)\sin\pi u}{3\pi}, \label{eq:Scos12} \\
  S^-_{\mathrm{sym}}(0,1;u)
    &= \frac{2(-2\sin\pi u + \sin 2\pi u)}{3\pi}, \label{eq:Ssin01} \\
  S^-_{\mathrm{sym}}(0,2;u)
    &= \frac{3\sin\pi u - \sin 3\pi u}{4\pi}. \label{eq:Ssin02}
\end{align}
The overlap differences $\Delta_j(u):=F_j(u)-G_j(u)$ satisfy:
\begin{enumerate}[label=\textup{(\alph*)}]
  \item For the first two energy slots, the individual differences are $O(1)$,
  but their sum satisfies
  $\Delta_{\mathrm{slot\,1}}(u) + \Delta_{\mathrm{slot\,2}}(u)
   = -(2+\sqrt{2})\,u + O(u^2)$.
  \item For $j\geq 2$: $\Delta_j(u)=O(u)$ individually.
\end{enumerate}
\end{lemma}

\begin{proof}
The profiles are definite integrals of smooth functions over
$[u-1,1]$, hence smooth in~$u$. For example,
$S^+_{\mathrm{sym}}(1,0;u)=\langle\varphi_1,\varphi_0(\cdot-uL)+\varphi_0(\cdot+uL)\rangle$
evaluates to $\frac{1}{L\sqrt{2}}\int_{(u-1)L}^{L}\cos(\pi t/L)\,dt + (\text{symmetric})
=\frac{\sqrt{2}\sin\pi u}{\pi}$
by direct integration; formulas~\eqref{eq:Scos12}--\eqref{eq:Ssin02} are obtained
similarly. At $u=0$, Fourier orthogonality gives $F_j(0)=G_j(0)=0$.
Differentiating~\eqref{eq:Ssin01}--\eqref{eq:Ssin02} at $u=0$: both sine
overlaps have zero derivative (Taylor expansions start at $u^2$).
For~\eqref{eq:Ssin01}: rewrite as
$-2\sin\pi u+\sin 2\pi u = 2\sin\pi u(\cos\pi u-1)$, vanishing to second order.
For~\eqref{eq:Ssin02}: rewrite as
$3\sin\pi u-\sin 3\pi u = 4\sin^3\pi u$, vanishing to third order.
The cosine derivatives are
$G_0'(0)=\sqrt{2}$ from~\eqref{eq:Scos10}
and $G_2'(0)=2$ from~\eqref{eq:Scos12}.
Therefore
$\sum_{\mathrm{slots}}\Delta_j'(0)=(0-\sqrt{2})+(0-2)=-(2+\sqrt{2})$.
Part~(b) follows from the mean value theorem applied to
$\Delta_j(u)=\Delta_j(u)-\Delta_j(0)$.
\end{proof}

\begin{lemma}[Higher-mode decay]
\label{lem:higher-mode}
Let $B_{1j}^\pm$ denote the coupling between the leading even/odd mode and
mode~$j\geq 2$ induced by the prime shifts, and let $g_j$ denote the spectral gap
between mode~$j$ and the leading mode in the Galerkin diagonal. Then:
\begin{enumerate}[label=(\roman*)]
  \item \emph{Coupling bound:}
  $|B_{1j}^\pm|\leq 4\pi\sqrt\lambda/L\cdot(1+O(1/L))$.
  \item \emph{Gap bound:}
  $g_j\geq 4\pi^2(j^2-1)\sqrt\lambda/L^2$ for all $j\geq 2$ and $L\geq 5$.
  \item \emph{Sector difference:}
  $|\Delta_{\mathrm{high}}|/D(\lambda)=O(1/L)\to 0$, where $D(\lambda)=\Theta(\sqrt\lambda)$
  is the three-mode frontier advantage.
\end{enumerate}
\end{lemma}

\begin{proof}
\emph{(i) Coupling bound.} For $j\neq 1$, orthogonality of the cosine basis on
$[-L,L]$ gives $\langle\varphi_1,(S_0+S_0)\varphi_j\rangle=0$. The mean-value
theorem yields $|\langle\varphi_1,(S_\delta+S_{-\delta})\varphi_j\rangle|\leq\pi\delta/L$.
Summing over primes with Abel summation~\cite{IwaniecKowalski2004} and
$\theta(x)=x+O(x\,e^{-c\sqrt{\log x}})$:
\[
  |B_{1j}^+|\leq\frac{\pi}{L}\sum_{p\leq\lambda}\frac{(\log p)^2}{\sqrt p}
  =\frac{4\pi\sqrt\lambda}{L}\bigl(1+O(1/L)\bigr).
\]
The same bound holds for $|B_{0j}^-|$ in the odd sector.

\emph{(ii) Gap bound.} By \Cref{def:auto-overlap}, $h_n(1)=(-1)^n/4$.
The prime sum concentrates at $p\sim\lambda$ (i.e., $u=\log p/\log\lambda\to 1$);
by Abel summation with the PNT~\cite{IwaniecKowalski2004}, the diagonal satisfies
$(W_{\mathrm{prime}})_{nn}\sim 4h_n(1)\sqrt\lambda=(-1)^n\sqrt\lambda$.
For even~$j$: $h_j(1)-h_1(1)=\tfrac12$, giving $g_j\geq\sqrt\lambda$.
For odd $j\geq 3$: $h_j(1)=h_1(1)=-\tfrac14$ (leading terms cancel).
Expanding $h_j(u)-h_1(u)$ about $u=1$ to second order gives
$h_j''(1)-h_1''(1)=\pi^2(j^2-1)/4$, whence
$g_j=4\pi^2(j^2-1)\sqrt\lambda/L^2+O(\sqrt\lambda/L^3)$.
In particular, $c_{\mathrm{gap}}:=4\pi^2/L\geq 2.8$ for $L\leq 14$ ($\lambda\leq 1{,}200{,}000$).

\emph{(iii) Sector difference.} Let
$E_{\cos}:=\sum_{j\geq 2}|B_{1j}^+|^2/g_j^+$ and
$E_{\sin}:=\sum_{j\geq 2}|B_{0j}^-|^2/g_j^-$
denote the resolvent-damped coupling energies in the even and odd sectors.
Each is $O(\sqrt\lambda)$ individually by~(i)--(ii). The sector
\emph{difference} $|E_{\sin}-E_{\cos}|$ is further suppressed by the Leading-Mode
Cancellation (\Cref{lem:leading-cancel}): the first-order Taylor coefficients of
the even and odd overlap differences at $u=0$ sum to $-(2+\sqrt{2})$, so the
leading contributions cancel pairwise. The coefficients
$c_{\sin}(L)-c_{\cos}(L)=O(1/L^2)$, and therefore
$|\Delta_{\mathrm{high}}|/D(\lambda)=O(1/L)\to 0$.
\end{proof}

\begin{remark}[Cancellation region for the sector difference]
\label{rem:leading-cancellation-region}
The cancellation used in part~(iii) of \Cref{lem:higher-mode} via
\Cref{lem:leading-cancel} is a Taylor identity at $u=0$, whereas the prime
sum that drives $D(\lambda)$ is dominated by frontier shifts with
$u=\log p/\log\lambda\in[0.7,1]$. At the Laplace endpoint $u=1$ the relevant
auto-overlap values are $h_n(1)=(-1)^n/4$ (cf.\ \Cref{def:auto-overlap}), so
the even and odd sectors are \emph{anti-correlated} rather than cancelling.
The $O(1/L)$-suppression claimed in part~(iii) therefore relies on the
averaging of the Taylor coefficient $c_{\cos}(L)-c_{\sin}(L)$ across the
slot structure described in \Cref{lem:leading-cancel}(a), not on a pointwise
cancellation at frontier shifts. A direct frontier-endpoint cancellation
identity in the spirit of \Cref{lem:leading-cancel} but anchored at $u=1$
would strengthen this estimate; in its absence, the bound should be regarded
as relying on the averaged Taylor structure of the slot decomposition and is
quantitatively conservative.
\end{remark}

\begin{lemma}[Galerkin truncation error]
\label{lem:galerkin}
The even sector is dominated by the frontier mode ($n=1$, diagonal
$\sim-\sqrt\lambda$); modes $n=0,2,3$ contribute at lower order, so $N_{\cos}=4$
suffices (\Cref{rem:galerkin-asymmetry}). The odd sector lacks this structural
amplification and is enclosed with $N_0=N_{\sin}\in[40,90]$. By the
Schur-complement formula~\cite{Kato1966}, the eigenvalue perturbation from
modes $j>N_0$ is bounded by
\[
  |\lambda_1^\pm(\mathrm{full})-\lambda_1^\pm(N_0)|
  \;\leq\;
  \frac{\norm{B_{\mathrm{tail}}}_{\mathrm{op}}^2}{g_{N_0+1}-\norm{B_{\mathrm{tail}}}_{\mathrm{op}}}.
\]
Substituting \Cref{lem:higher-mode}(i)--(ii) gives
$\norm{B_{\mathrm{tail}}}_{\mathrm{op}}\leq 8\pi\sqrt\lambda/L$ and
$g_{N_0+1}\geq 4\pi^2 N_0^2\sqrt\lambda/L^2$, so
\[
  |\lambda_1^\pm(\mathrm{full})-\lambda_1^\pm(N_0)|
  \;\leq\;
  \frac{16\,\sqrt\lambda}{N_0^2}\,(1+o(1))
  \quad\text{as }\lambda\to\infty.
\]
At every sampled grid point this analytical bound is dominated by the
reported Landscape safety factor of \Cref{rem:truncation-safety}
($\geq 34\times$ at $\lambda=100$, $\geq 65\times$ for $\lambda\geq 200$),
which is a consistency check for the reported gap, not a lower-bound proof.
\end{lemma}

\begin{proof}
The Schur-complement bound is standard~\cite{Kato1966}.
\Cref{lem:higher-mode}(i) gives $\norm{B_{\mathrm{tail}}}_{\mathrm{op}}\leq 8\pi\sqrt\lambda/L$.
\Cref{lem:higher-mode}(ii) gives $g_{N_0+1}\geq 4\pi^2 N_0^2\sqrt\lambda/L^2 \gg \norm{B_{\mathrm{tail}}}_{\mathrm{op}}$
for $N_0\geq 40$, hence the denominator is dominated by $g_{N_0+1}$.
The displayed bound $16\sqrt\lambda/N_0^2$ follows. The reported gap uses the
sharper interval-arithmetic enclosure of \cite[Section~2.9]{GeigerLandscapeII}
with the diagnostic Cauchy tail bound described in \Cref{rem:truncation-safety}.
\end{proof}

\section{Main Theorem and Open Conjecture}

\paragraph{v1.5 status of the main theorem.}  In the Connes--van
Suijlekom form of Section~2, the main theorem statement (RH
conditional on the Connes program and the Asymptotic Variational Gap
Conjecture) is unchanged.  What changed is the precise formulation of
the conditioning:
\begin{itemize}[leftmargin=*]
  \item \emph{External preprint input from Connes 2026 §6.5
    (Fact 6.4):} the Fourier transform of the educated-guess
    eigenfunction $k_\Lambda$ converges to $\Xi$ uniformly on substrips
    $|\Im z|<\delta<1/2$ at rate $O(\Lambda^{-1/2-\delta})$.  This is
    stated in \cite{Connes2026} as a consequence of the
    classical Slepian--Pollak prolate spheroidal wave function theory
    ([Connes2026, Theorem~1 of reference~47]), not as a conjecture.
  \item \emph{External preprint input from Connes--van Suijlekom 2025
    \cite{ConnesvS2025}:} the fixed-$\Lambda$ quadratic-form
    real-zero criterion, proved within that preprint.
  \item \emph{Internal open question:} simplicity and parity of the
    smallest eigenvalue of $QW_\Lambda$ (Connes 2026 §6.6 ``remaining
    steps'' item~(i)) --- this is exactly even dominance as defined in
    Definition~\ref{def:even-dominance}, addressed in the present paper.
  \item \emph{Internal open question:} approximation quality
    $k_\Lambda\approx\xi_\Lambda$ (Connes 2026 §6.6 ``remaining steps''
    item~(ii)) --- this is the Mass-Concentration question MS2
    addressed in the companion Zookeeper paper \cite{GeigerZookeeper2026}.
\end{itemize}
The conditioning in v1.5 is therefore that
\emph{even dominance for $QW_\Lambda$} holds along a sequence
$\Lambda_n\to\infty$, together with the Connes--van Suijlekom
preprint criterion, gives RH via the Fact~6.4 bridge stated by
Connes~2026.  The companion
Zookeeper question is independent and concerns whether the specific
educated-guess approximant produces additional structural control.

\subsection{Certified Range: Finite Even Dominance}

\begin{theorem}[Finite even dominance under lower-bound validation]
\label{thm:even-finite}
Assume that the finite certificate matrices are recomputed in the corrected
Connes--van Suijlekom framework and that the quantities reported as
$\underline{\lambda_1^-}(\lambda)$ in \Cref{prop:finite-certificates} are
validated as rigorous odd-sector lower bounds at each grid point.
Then for every $\lambda$ in that grid (33 values spanning
$100\leq\lambda\leq 1{,}300{,}000$), $\QW_\lambda$ has a simple even ground state,
i.e.\ $\lambda_1^+(\lambda)<\lambda_1^-(\lambda)$ and the even minimum is
isolated.
\end{theorem}

\begin{proof}
Under the additional lower-bound validation hypothesis, the interval-arithmetic
Galerkin diagnostics of \Cref{prop:finite-certificates} produce, at each
$\lambda$ in the grid, a rigorous upper bound
$\overline{\lambda_1^+}(\lambda)$ and a rigorous lower bound
$\underline{\lambda_1^-}(\lambda)$ from the even and odd Galerkin blocks
respectively, together with validated tail bounds. The resulting gap
$\overline{\lambda_1^+}(\lambda)-\underline{\lambda_1^-}(\lambda)<0$ at every
grid point directly establishes $\lambda_1^+(\lambda)<\lambda_1^-(\lambda)$ on
the grid. Simplicity follows from the validated intra-even gap
$\lambda_2^+-\lambda_1^+>0$ at each grid point.
\end{proof}

\subsection{Asymptotic Regime: Variational Statement and Open Conjecture}

The asymptotic argument of \Cref{prop:frontier-dominance} establishes the
\emph{variational} statement
\begin{equation}\label{eq:variational-result}
  \lambda_{\min}\bigl(W^+_{\mathrm{prime}}(\lambda)-W^-_{\mathrm{prime}}(\lambda)\bigr)
  \;=\; -\Theta(\sqrt\lambda)
\end{equation}
on the three-mode Galerkin block, with explicit threshold
$\Lambda_*\leq 175{,}000$ (\Cref{rem:lambda-star-improved}).
This is the statement that the difference $W^+-W^-$ has a negative minimum
eigenvalue with quantitative rate of growth.

\paragraph{The variational direction gap.}
Statement~\eqref{eq:variational-result} is \emph{not} the same as
$\lambda_1^+(\lambda)<\lambda_1^-(\lambda)$. The Rayleigh inequality
$\lambda_{\min}(W^+-W^-)\leq v^T(W^+-W^-)v$ gives an \emph{upper} bound on the
minimum of the difference, equivalently it shows the existence of a vector $v$
with $v^TW^+v<v^TW^-v$. For even dominance one needs, in addition, a
quantitative \emph{lower} bound on $\lambda_1^-(\lambda)$ that places the odd
minimum above the certified upper bound for $\lambda_1^+$. The simple
$2\times2$ example
\[
  A=\begin{pmatrix}0&0\\0&10\end{pmatrix},\quad
  B=\begin{pmatrix}5&0\\0&-5\end{pmatrix}
\]
shows that $\lambda_{\min}(A-B)=-5<0$ does not imply
$\lambda_{\min}(A)<\lambda_{\min}(B)$; in fact $\lambda_{\min}(A)=0>-5=\lambda_{\min}(B)$.
The finite-range data support this picture but do not close the gap at
theorem level until the odd-sector lower bound is validated; asymptotically a
separate argument is required.

\begin{conjecture}[Asymptotic Variational Gap Conjecture]
\label{conj:avg}
There exist constants $\Lambda^\dagger>0$, $c_+>c_->0$ and an exponent
$\beta\in(0,1)$ such that for all $\lambda\geq\Lambda^\dagger$,
\[
  \lambda_1^+(\lambda) \;\leq\; -\,c_+\sqrt\lambda,
  \qquad
  \lambda_1^-(\lambda) \;\geq\; -\,c_+\sqrt\lambda \,+\, c_-\sqrt\lambda\,L^{-\beta},
\]
where $L=\log\lambda$. In particular, $\lambda_1^+(\lambda)<\lambda_1^-(\lambda)$
for all $\lambda\geq\Lambda^\dagger$.
\end{conjecture}

\Cref{sec:asymptotic-lower-bound} discusses what would be required to establish
\Cref{conj:avg}, sketches a proposed approach via degenerate perturbation
theory on the asymmetric three-mode diagonal structure, and explains why the
naive variational argument cannot be lifted directly.

\subsection{Riemann Hypothesis: Conditional Statement}

\begin{theorem}[Riemann Hypothesis, conditional]
\label{thm:rh}
Assume \Cref{conj:avg} (Asymptotic Variational Gap) and the Connes program,
namely the quadratic-form real-zero criterion of~\cite{ConnesvS2025} together
with the Hurwitz bridge of \cite[Fact~6.4 and Section~6.4--6.6]{Connes2026}.
Then all non-trivial zeros of the Riemann zeta function~\cite{Riemann1859}
lie on the critical line $\Re(s)=1/2$.
\end{theorem}

\begin{proof}
By \Cref{conj:avg}, even dominance holds for every
$\lambda\geq\Lambda^\dagger$, hence along a divergent sequence
$\lambda_n\to\infty$. If the finite certificate lower bounds are later validated,
\Cref{thm:even-finite} adds a finite-range bridge, but it is not needed for this
divergent sequence. Applying \Cref{cor:sequence-suffices} to the sequence: the Fourier
transforms of the even-dominant Weil forms have only real zeros
(\Cref{thm:connes}), and the Hurwitz bridge transfers the real-zero property
to the limiting $\Xi$-function. Therefore the non-trivial zeros of $\zeta(s)$
lie on $\Re(s)=1/2$.
\end{proof}

\paragraph{Remark on conditioning.}
\Cref{thm:rh} carries two distinct conditional inputs. The Connes--van
Suijlekom criterion and the Hurwitz bridge are external results of the Connes
program~\cite{Connes2026,ConnesvS2025}, currently arXiv preprints; their
analytical structure is discussed in \Cref{rem:hurwitz-analytical}. The
Asymptotic Variational Gap Conjecture~\ref{conj:avg} is internal to the
present approach and is the central open mathematical problem for closing
the asymptotic half of even dominance unconditionally within this framework.

\section{Status of the Asymptotic Lower Bound for \texorpdfstring{$\lambda_1^-$}{lambda\_1\^{}-}}
\label{sec:asymptotic-lower-bound}

This section is not part of the proof of \Cref{thm:rh}: it explains why the
common-Rayleigh-vector argument of \Cref{prop:frontier-dominance} cannot be
lifted directly to even dominance in the asymptotic regime, and outlines a
proposed approach -- via degenerate perturbation theory on the asymmetric
three-mode diagonal structure -- toward establishing \Cref{conj:avg}.

\subsection{Why the Variational Bound Is Not Even Dominance}

\Cref{prop:frontier-dominance} establishes
$\lambda_{\min}(D_{\mathrm{tot}}(\lambda))<0$ on the three-mode block, where
$D_{\mathrm{tot}}=W^+_{\mathrm{prime}}-W^-_{\mathrm{prime}}$ (restricted to the
$3\times3$ block). By the Rayleigh variational inequality this means: there
exists a vector $v$ (namely $v=e_1$) with $v^TW^+v<v^TW^-v$. It does not
follow that the smallest eigenvalue of $W^+$ lies below the smallest eigenvalue
of $W^-$. Indeed, $\lambda_{\min}(W^\pm)\leq v^TW^\pm v$ holds for any unit
vector $v$, so both bounds run in the same direction:
$\lambda_{\min}(W^+)\leq e_1^TW^+e_1$ and $\lambda_{\min}(W^-)\leq e_1^TW^-e_1$.
A quantitative \emph{lower} bound for $\lambda_{\min}(W^-)$ is needed.

\subsection{Asymmetric Three-Mode Diagonal Structure}

At the Laplace endpoint $u=1$ (the frontier regime
$r=\log p/\log\lambda\approx 1$, which dominates the prime sum), the asymptotic
leading-order diagonal entries of $W^\pm_{\mathrm{prime}}$ on the three-mode
block are
\[
  \mathrm{diag}\bigl(W^+_{\mathrm{prime}}\bigr)
  \;\sim\; \sqrt\lambda\cdot(+1,\,-1,\,+1),
  \qquad
  \mathrm{diag}\bigl(W^-_{\mathrm{prime}}\bigr)
  \;\sim\; \sqrt\lambda\cdot(-1,\,+1,\,-1).
\]
These signs follow from the auto-overlap profile at $u=1$:
$h_n(1)=(-1)^n/4$ for even basis functions and $h_n^{\sin}(1)=(-1)^{n+1}/4$
for odd, weighted by the frontier prime sum
$\sum_{p\in\text{frontier}}\log p/\sqrt p\sim 2\sqrt\lambda$.

The structural asymmetry is then visible:
\begin{itemize}[leftmargin=*]
  \item In $W^+$ the minimum diagonal $(-\sqrt\lambda)$ is at index $1$ and is
  \emph{isolated}: the other two diagonals are $(+\sqrt\lambda)$, an
  $O(\sqrt\lambda)$-distance away. Off-diagonal couplings shift $\lambda_1^+$
  by only \emph{second-order} amounts.
  \item In $W^-$ the minimum diagonals $(-\sqrt\lambda)$ are at indices $0$ and
  $2$ and are \emph{near-degenerate}, with the central diagonal $(+\sqrt\lambda)$
  separated by $O(\sqrt\lambda)$. Off-diagonal couplings between the degenerate
  pair split them in \emph{first} order, potentially driving $\lambda_1^-$
  further down.
\end{itemize}

A naive analysis stops here and yields \emph{anti-}even-dominance:
$\lambda_1^-<\lambda_1^+$, contradicting the 33 finite-range certificates. The
resolution must come from a subleading order of the relevant off-diagonal
coupling $b_{02}:=(W^-_{\mathrm{prime}})_{02}$ at the Laplace endpoint.

\subsection{Elementary Cancellation in the Degenerate Odd Coupling}

The off-diagonal coupling $b_{02}$ in the odd sector at shift $r$ involves
the overlap integral
\[
  I_-(r) \;=\; \int_0^1 \sin(\pi x)\sin\bigl(3\pi(x-r)\bigr)\,dx.
\]
At $r=1$ this vanishes by elementary sine orthogonality on $[0,1]$:
$\sin(\theta-3\pi)=\sin(\theta-\pi)=-\sin(\theta)$, so
\[
  I_-(1) \;=\; -\int_0^1 \sin(\pi x)\sin(3\pi x)\,dx \;=\; 0.
\]
Differentiability of $I_-$ at $r=1$ gives $I_-(r)=O(1-r)$ near the endpoint.
In the frontier window $r\in[0.7,1]$, the weighted prime sum
$\sum_p (\log p/\sqrt p)\cdot I_-(r_p)$ with $\langle 1-r_p\rangle\sim 1/L$
on average should therefore satisfy
\[
  \bigl|b_{02}(\lambda)\bigr| \;=\; O\bigl(\sqrt\lambda/L\bigr).
\]
If this estimate can be made rigorous with explicit constants, the degenerate
splitting in $\lambda_1^-$ is at most $O(\sqrt\lambda/L)$, which is
\emph{subleading} compared to a possible $-c_+\sqrt\lambda$ shift in
$\lambda_1^+$ from the same frontier mechanism.

\subsection{Proposed Three-Step Programme}

\Cref{conj:avg} would follow from the following programme:

\paragraph{Step 1. Diagonal accounting.}
Compute the explicit diagonal entries of $W^\pm_{\mathrm{prime}}$ on the three-mode block, separating leading
frontier contributions from subleading tail and archimedean terms. This is
bookkeeping without new mathematics, and should produce explicit constants
$c^\pm_{nn}$ such that $(W^\pm_{\mathrm{prime}})_{nn}=c^\pm_{nn}\sqrt\lambda+o(\sqrt\lambda)$.

\paragraph{Step 2 (Off-diagonal control).} Compute the trigonometric overlap
integrals $I_{nm}^\pm(r)=\int_0^1 \varphi_n^\pm(x)\varphi_m^\pm(x-r)\,dx$ for
$n,m\in\{0,1,2\}$ and $r\in[0.7,1]$, and aggregate against the frontier prime
sum. Establish the quantitative bound
$|b_{02}(\lambda)|\leq C\sqrt\lambda/L$ with explicit $C$.

\paragraph{Step 3 (Degenerate perturbation theory).} Apply Bauer--Fike or a
refined perturbation theorem for near-degenerate eigenvalues with controlled
coupling. Obtain
\[
  \lambda_1^-(\lambda) \;\geq\; -\sqrt\lambda\,(1 + O(1/L)),
\]
combine with the upper bound
$\lambda_1^+(\lambda)\leq -\sqrt\lambda\,(1 + c_+/\text{slowly-varying})$ from
the common-Rayleigh argument (after rederivation in absolute terms, not
through the difference), and conclude that
$\lambda_1^+<\lambda_1^-$ asymptotically.

\paragraph{Step 4 (Galerkin transfer).} Use the existing Schur-complement
machinery (\Cref{lem:higher-mode} and \Cref{lem:galerkin}) to transfer the
three-mode statement to the full Weil operator $A_\lambda^\pm$. This is
already in place modulo the issue noted in
\Cref{rem:leading-cancellation-region} below.

\paragraph{Honest assessment.} Steps~1 and~2 are routine but laborious
(several weeks). Step~3 is genuinely difficult: degenerate perturbation theory
with sign-controlled couplings is an established discipline, but every
explicit constant can become a subproblem. The relationship
$|\lambda_1^+-\lambda_1^-|=O(\sqrt\lambda/L)$ that this programme suggests is
consistent with the numerical observation that the reported gap is roughly
$20\%$ of $|\lambda_1^+|$ -- not a constant fraction of $\sqrt\lambda$, but a
logarithmically shrinking fraction.

\subsection{Update (14 May 2026): Strengthened Structural Concern after
            Step-1 Bookkeeping}
\label{sec:k0-strengthened}

A careful execution of Step~1 above (separating leading frontier
diagonals via Laplace asymptotics, with numerical validation) has
produced a stronger structural concern that the reviewer's original
sketch did not anticipate. We record it honestly:

\paragraph{Symmetric-shift diagonal equivalence.} In the basis of
Section~\ref{sec:parity-matrix}, the symmetric shift $S_u + S_{-u}$
applied to the cosine basis $\varphi_n(t) = \cos(n\pi t/L)/\sqrt L$
($n\geq 1$) and to the sine basis $\psi_{n-1}(t) = \sin(n\pi t/L)/\sqrt L$
gives \emph{identical} auto-overlap profiles
\[
  \langle\varphi_n, (S_u+S_{-u})\varphi_n\rangle
  \;=\; \langle\psi_{n-1}, (S_u+S_{-u})\psi_{n-1}\rangle
  \;=\; 2\,(1 - u/(2L))\,\cos(n\pi u/L),
\]
for $u \in [0,2L]$. The sine-correction terms in the asymmetric
shifts cancel under symmetrization. Numerical verification
(script \texttt{verify\_diagonal\_fast.py}) confirms agreement of the
corresponding prime-summed diagonal entries to machine precision for
$\lambda \in \{100, 300, 1000\}$.

\paragraph{Spectral consequence.} Since
$\langle\varphi_n, (S_u+S_{-u})\varphi_m\rangle = 0$ for $n\neq m$ by
Fourier orthogonality (the sine-cross-terms cancel under symmetrization),
the Weil operator $A_\lambda$ is exactly diagonal in this Galerkin
basis, with diagonal entries
\[
  \mu^+_n
  = \langle\varphi_n, A_\lambda\,\varphi_n\rangle,
  \quad
  \mu^-_{n-1}
  = \langle\psi_{n-1}, A_\lambda\,\psi_{n-1}\rangle.
\]
Combined with the auto-overlap equivalence, this gives
$\mu^+_n = \mu^-_{n-1}$ for every $n \geq 1$, both in the prime
contribution and (by the same symmetric-shift argument) in the
archimedean contribution. Therefore in this Galerkin model,
\[
  \sigma(A_\lambda^+) \;=\; \sigma(A_\lambda^-)\,\cup\,\{\mu^+_0\},
\]
where $\mu^+_0$ is the constant-mode eigenvalue. Numerically
$\mu^+_0 > 0$ and large, hence not the minimum of $\sigma(A_\lambda^+)$.
The conclusion is
\[
  \lambda_1^+(\lambda) \;=\; \lambda_1^-(\lambda) \quad\text{exactly}
\]
in this Galerkin model, i.e., even dominance is \emph{not} attained
asymptotically.

\paragraph{Tension with the finite-range certificates.}
The 33 IA-certificates of \cref{prop:finite-certificates} report
$\text{cert\_gap} := \overline{\lambda_1^+}_{\rm CAP} - \underline{\lambda_1^-}_{\rm CAP} < 0$,
which is incompatible with the above structural equivalence unless
either the spectral-equivalence statement is wrong or the
$\underline{\lambda_1^-}_{\rm CAP}$ is not a rigorous lower bound. An
audit of the certifier code \texttt{certifier\_production.py} shows
that $\underline{\lambda_1^-}_{\rm CAP}$ is computed as
$\lambda_1^{\rm Galerkin}(N_{\rm big}=60) - \text{remaining} - \epsilon$,
where $\text{remaining}$ is a heuristic geometric extrapolation of
column norms (no Lehmann--Maehly, Kato--Temple, or rigorous
Schur--complement bound is applied). The Galerkin variational
principle gives $\lambda_1^{\rm Galerkin}$ as an \emph{upper} bound on
$\lambda_1^-$, not a lower bound, and subtracting a heuristic
correction does not yield a rigorous lower bound without further
analysis. The certificates should therefore be read as
\emph{computer-assisted numerical evidence}, not as rigorous
mathematical proof, until the lower-bound construction is replaced by
a method backed by a quantitative perturbation theorem.

\paragraph{Honest current status.} Both pillars of the even-dominance
argument (asymptotic frontier dominance and finite-range CAP
certificates) are under structural concern as of this revision. We
do not at present claim even dominance for $\QW_\lambda$, neither
asymptotically nor on the finite diagnostic range, unless one of the
following can be established:
\begin{itemize}[leftmargin=*]
  \item the Galerkin operator in \cref{eq:quad-form-connesvs} differs from the
  actual Connes operator in a way that breaks the cosine--sine
  frequency symmetry; or
  \item the heuristic tail correction in
  $\underline{\lambda_1^-}_{\rm CAP}$ can be replaced by a quantitative
  perturbation-theoretic lower bound, and the resulting rigorous
  reported \texttt{cert\_gap} is still negative.
\end{itemize}
The Shift Parity Lemma (\cref{thm:shift-parity}), the explicit
Laplace-endpoint diagonal formulas, the non-existence theorems NE-A
and NE-B (Landscape Part~II), and the numerical phenomenology of the
reported gap are unaffected by this concern and remain valid as
independent structural results.

This update is in the spirit of treating the present manuscript as a
continuously evolving draft. The companion notes
\path{_proof-notes/K0_STEP1_FINDINGS.md} and
\path{_proof-notes/K0_CAP_AUDIT.md} provide the full analytical
and numerical record of the present concern.

\subsection{Update (15 May 2026): Connes-vS Operator Form Resolves the
            Structural Tension; Rigorous CAP Closure Remains Open}
\label{sec:k0-cvs-resolution}

A subsequent comparison of the operator definition against
Connes--van Suijlekom \cite{ConnesvS2025} clarified that the strengthened
K0 concern (the spectral equivalence
$\sigma(A^+) = \sigma(A^-) \cup \{\mu^+_0\}$) was an artefact of the
legacy operator realisation on $L^2([-L,L])$ with real cosine/sine
basis under $t\mapsto -t$ symmetry, not of the actual Connes operator.
The proven Connes--van Suijlekom operator acts on $L^2([0,L])$ with
the complex exponential basis $U_n$, and the cosine/sine parity
sectors are defined via the frequency reflection
$U_n\leftrightarrow U_{-n}$.  In this corrected formulation
(Section~2 v1.5 reformulation) the diagonals
$\langle\cos_n|Q|\cos_n\rangle = D_n - S_n$ and
$\langle\sin_n|Q|\sin_n\rangle = D_n + S_n$ differ by the genuinely
non-trivial asymmetry $2S_n$, where
$S_n = (\pi n)^{-1}\int_0^L\sin(2\pi n y/L)\,dD(y)$.  Direct
computation (script \path{_proof-notes/verify_connes_vs_form.py})
gives $S_n > 0$ for the Weil distribution and all tested $\lambda$.
For $\lambda=1000$, $n=1$: $S_1\approx 8.63$, cos diagonal
$\approx 3.26$, sin diagonal $\approx 20.52$.  Asymptotically
$S_1\sim 8\sqrt\lambda/\log\lambda$ (script
\path{verify_connes_vs_form.py}, derivation in
\path{K0_FRONTIER_SN_ASYMPTOTIC.md}). The structural foundation
for even dominance is therefore present in the correct operator form.

The remaining task is a \emph{rigorous} computer-assisted-proof closure
of $\lambda_1^{\cos}(\lambda) < \lambda_1^{\sin}(\lambda)$ in this
form.  A systematic analysis of the available rigorous lower-bound
techniques is recorded in
\path{_proof-notes/K0_LEHMANN_ANALYSIS.md}, with the conclusion:

\begin{itemize}[leftmargin=*]
  \item Temple's inequality
    $\lambda_1\geq\mu-\eta^2/(\rho-\mu)$ requires
    $\rho\leq\lambda_2$ — an \emph{a priori lower bound} on the second
    eigenvalue, which the Galerkin variational principle does not
    provide.  The naive substitution $\rho=\mu_2$
    \emph{is not rigorous}: $\mu_2\geq\lambda_2$ by Cauchy
    interlacing, hence using $\mu_2$ in the denominator underestimates
    the correction.

  \item Lehmann--Maehly's generalised inequality requires
    $\sigma\leq\lambda_{N+1}$, again an a priori lower bound on a
    higher eigenvalue.

  \item Bauer--Fike identifies the \emph{nearest} eigenvalue to $\mu$,
    not specifically the lowest, and so is ambiguous without
    additional gap information.

  \item Operator-norm bounds give $\lambda_1\geq-\|A\|_{\rm op}$ with
    $\|A_\lambda\|_{\rm op}\le 4\sqrt\lambda + O(\log\lambda)$, which
    is far too loose.

  \item Schur-complement and Gerschgorin bounds require either a
    lower bound on the complementary block or diagonal dominance.
    Gerschgorin fails for the Connes--van Suijlekom Galerkin matrix
    because the off-diagonal row sums diverge logarithmically: the
    cross-overlap $\sin/(\pi(n-m))$ has only $1/|n-m|$ decay.
\end{itemize}

The common underlying obstacle is the absence of an
\emph{a priori lower bound on $\lambda_2$} or on the tail spectrum,
which Standard CAP techniques uniformly require.  Three medium-term
research directions to address this are sketched in
\texttt{K0\_LEHMANN\_ANALYSIS.md}: (i) Mollification of the Weil
distribution to discrete spectrum, followed by stability of the
$\epsilon\to 0$ limit; (ii) Toeplitz approximation and symbol-based
perturbation theory exploiting the near-Toeplitz structure of the
Connes--van Suijlekom Galerkin matrix; (iii) reverse-engineering the
proof technique of Connes--vS 2025 (which contains an internal Hurwitz
argument and thus must implicitly handle related lower-bound
questions).

\paragraph{Current honest position.} The Connes--van Suijlekom operator
form is the correct setting and exhibits even dominance numerically
with consistent margin across $\lambda$ and $N$.  The asymptotic
$S_1$-asymptotic gives a quantitative Cos--Sin diagonal gap of
order $\sqrt\lambda/\log\lambda$.  A rigorous CAP-style closure in the
sense of producing a certified lower bound on $\lambda_1^{\sin}$ that
strictly exceeds a certified upper bound on $\lambda_1^{\cos}$ remains
open and requires a non-trivial technical advance beyond
straightforward Galerkin + Temple/Lehmann.  We document this honestly
in line with the convention that the present manuscript is a
continuously developing draft.

% added 2026-05-23: Route B parallel development reference
\subsection{Update (May 2026): Parallel Route via Connes-vS Err-Collapse and I-1 Normal-Family Reframe}
\label{sec:route-b-parallel}

A parallel conditional closure of the Riemann Hypothesis via a
different handling of wall~(ii) has been developed in the companion
proof notes (see \path{_proof-notes/K0_NORMAL_FAMILY_REFRAME_2026-05-22.md}).
Recall that the present proof delegates wall~(ii)
(the approximation-quality condition $k_\Lambda \approx \xi_\Lambda$,
i.e.\ the MS2 bound) to the companion Zookeeper paper.
The parallel route --- referred to as Route~B~--- shows that the same
spectral structure gives a second conditional closure via a
normal-family (Montel--Vitali) argument: wall~(ii) decomposes
rigorously into an a-priori strip bound~(ii-a) on the family
$\{\hat\xi_\lambda\}$ (the canonical entire functions of Paley--Wiener
type $\hat\xi(z) = (2/\sqrt L)\sin(zL/2)\cdot F(z)$, centered, even,
real-on-real) and a limit-identification condition~(ii-c); and (ii-c)
is absorbed into~(ii-a) rigorously via parity and Hadamard order
(modulo a zero-coincidence assumption~(Z), \S{}E.14 of the notes).
The remaining open blocker for Route~B is therefore the single
a-priori uniform strip bound~(ii-a) from $QW_N$ coercivity.

Converged empirical measurements on the canonical objects (Mac Studio,
$\mathrm{dps} \approx 30\text{--}50\cdot\lambda$, $\lambda \in
\{3,5,7,9,11,13,15\}$, \S{}E.18 of the notes; values at $\lambda=11,13$
bit-identical across two independent dps-runs, and $\lambda=15$
bit-identical at $\mathrm{dps}=640$ and $\mathrm{dps}=800$) yield
$A_\lambda(0.4) \approx 1.003\text{--}1.004$ \emph{flat} across the
seven converged $\lambda$-values
(the SUP-observable
$A_\lambda(\delta) := \sup_{x\in\mathbb{R}} |\hat\xi_\lambda(x)|^2 /
(1+|x|^{2\delta})$, normalised) and $B_\lambda(0.4)/\lambda^{0.4}
\approx 0.68$ plateau (the $L^2$-observable, sub-generic growth):

\begin{center}
\footnotesize
\renewcommand{\arraystretch}{1.15}
\begin{tabular}{c|c|c|c|c}
$\lambda$ & $\sup|\hat\xi|_{\mathbb{R}}$ & $A_\lambda(0.4)$ & $B_\lambda(0.4)$ & $B_\lambda/\lambda^{0.4}$ \\\hline
 3 & 0.858 & 1.0033 & 1.151 & 0.74 \\
 5 & 0.879 & 1.0037 & 1.340 & 0.70 \\
 7 & 0.886 & 1.0038 & 1.505 & 0.69 \\
 9 & 0.890 & 1.0039 & 1.649 & 0.68 \\
11 & 0.890 & 1.0039 & 1.778 & 0.68 \\
13 & 0.891 & 1.0039 & 1.895 & 0.68 \\
15 & 0.892 & 1.00394 & 2.003 & 0.68
\end{tabular}
\end{center}

The flatness of $A_\lambda$ indicates that the spectral framework is
well-conditioned and consistent with a uniform strip bound; it does not
constitute a proof of~(ii-a).  Earlier measurements reported in
\S{}E.16--E.17 of the notes were retracted due to insufficient
precision ($\mathrm{dps}=30$) near cluster near-degeneracy of the
minimal even eigenvalue; only the converged \S{}E.18 values are
reliable.

\section{What Is Not Used}

The proof does not use the following Landscape material:
\begin{itemize}[leftmargin=*]
  \item the original gauge-thermodynamic route and its obstructions;
  \item an absolute $\mathrm{PF}_\infty$ or total-positivity claim for the full prime operator;
  \item a universal commuting Sturm--Liouville operator;
  \item Hellmann--Feynman monotonicity between certificate points;
  \item the narrative Dungeon/Memorial presentation of excluded routes.
\end{itemize}
Those results remain valuable as route classification and audit history, but
they are not part of the proof chain above.

\appendix

\section{Constants Audit Table}
\label{sec:constants-audit}

All explicit numerical constants used in the proof, with their source, role, and
verification mechanism. Constants flagged ``IA'' are interval-arithmetic certified
(\texttt{mpmath.iv}, $50$-digit precision); ``CAP'' are the $33$ Landscape
computer-assisted certificates; ``PNT/Mertens'' are explicit asymptotic remainders.

\begin{center}
\footnotesize
\renewcommand{\arraystretch}{1.25}
\begin{tabular}{p{2.6cm}|p{2.0cm}|p{6.2cm}|p{2.6cm}}
\textbf{Constant} & \textbf{Value} & \textbf{Role / Source} & \textbf{Verification} \\\hline
$C_f$ & $0.4685$ & Common frontier Rayleigh quotient $-d_{11}(r)\geq C_f$ on $r\in[0.7,1]$ (\Cref{lem:common-vector}) & IA on $[0.7,1]$ (1D)\\
$C_n$ & $2.2847$ & Operator-norm bound $\sup_{r\in(0,2)}\norm{D_3(r)}_{\mathrm{op}}\leq C_n$ (\Cref{lem:nonfrontier}) & IA on $(0,2)$ (3D)\\
$c_{\mathrm{gap}}$ & $\geq 4\pi^2/L\geq 2.8$ & Spectral-gap constant for $L\leq 14$ (\Cref{lem:higher-mode}(ii)) & Laplace asymptotics\\
$N_{\cos}$ & $4$ & Even-block Galerkin truncation (\Cref{rem:galerkin-asymmetry}) & Stability check $k=3,4,5$\\
$N_{\sin}$ & $40$--$90$ & Odd-block Galerkin truncation (\Cref{rem:galerkin-asymmetry}) & Cauchy tail bound\\
$\Lambda_*$ & $\leq 175{,}000$ & Asymptotic threshold; prime-sum computation (\Cref{rem:lambda-star-improved}); conservative bound $\leq 1{,}300{,}000$ (\Cref{prop:frontier-dominance}) & Exact prime enum.\\
$\theta(x)-x$ & $\leq 0.006788\,\tfrac{x}{\log x}$ & Effective Chebyshev bound for $x\geq 89{,}967{,}803$ (Dusart 2010, \cite{Dusart2010}) & Unconditional\\
$|\mathrm{cert\_gap}|$ at $\lambda{=}100$ & $\geq 2.25$ & First CAP (\Cref{prop:finite-certificates}) & CAP\\
$|\mathrm{cert\_gap}|$ at $\lambda{=}1.3{\cdot}10^6$ & $\geq 475.83$ & Last CAP & CAP\\
Safety factor at $\lambda{=}100$ & $\geq 34\times$ & $|\mathrm{cert\_gap}|$ vs.\ Galerkin/tail error (\Cref{rem:truncation-safety}) & IA + Cauchy tail\\
Safety factor at $\lambda\geq 200$ & $\geq 65\times$ & Same, monotone in $\lambda$ & IA + Cauchy tail\\
$c$ (slot cancellation) & $2+\sqrt 2$ & Closed-form constant (\Cref{lem:leading-cancel}(a)) & Symbolic (sympy)
\end{tabular}
\end{center}

The Landscape companion repository archives the verification scripts. The most
relevant entries are:
\begin{itemize}[leftmargin=*,itemsep=2pt]
  \item \texttt{shift\_parity\_cert\_v3\_targeted.py} -- Shift Parity Lemma;
  \item \texttt{certifier\_production.py} -- the $33$ CAP certificates;
  \item \texttt{partA\_proof\_sketch.py} -- frontier-dominance check;
  \item \texttt{euler\_maclaurin\_certifier.py} -- asymptotic threshold;
\end{itemize}
see~\cite[\path{scripts/}]{GeigerLandscapeII}. A dedicated companion repository
for the present extraction is mirrored at \projectrepo.

\section*{AI Disclosure}

\noindent\textbf{Statement on the Use of AI Tools}

\medskip

\begin{sloppypar}
\noindent
This manuscript was prepared with moderate assistance from AI-based language
models (Claude, Anthropic; Gemini, Google; ChatGPT, OpenAI). The use
encompassed the following areas:
\end{sloppypar}

\begin{itemize}
  \item \textit{Research:} Assistance with systematic literature search,
        identification and summarisation of relevant research literature.
  \item \textit{Structuring:} Support in outlining and logically organising
        the manuscript, in particular the proof-only extraction from the
        Landscape series.
  \item \textit{Drafting:} Assistance in formulating individual sections
        and paragraphs based on core arguments and theses specified by
        the author.
  \item \textit{Language:} Grammar and spell checking, linguistic revision.
  \item \textit{Formatting:} Preparation and verification of citations and
        bibliographies; \LaTeX{} typesetting support.
  \item \textit{Internal review:} A 7-phase model-based review chain
        (constructive, expert, adversarial) was used to identify and
        address structural and citation issues prior to submission.
\end{itemize}

\noindent
The entire research process -- including the research question, methodological
decisions, substantive argumentation, critical evaluation of sources, and
conclusions -- was conceived, directed, and borne by the author. All
computer-assisted certificates were computed by deterministic numerical
scripts independently of AI. All AI-generated suggestions were critically
reviewed, revised, and where appropriate rejected. The author reviewed all
content to the best of his knowledge but cannot guarantee complete
independent verification of all factual claims and references.

\medskip

\noindent
Full scholarly and substantive responsibility for this article rests with
the author, Lukas Geiger.

\bibliographystyle{plain}
\bibliography{references}

\end{document}
